
\Data\ID{1003057901}
\Updated{2020-07-15 03:07:12}
\Level{A}
\SubArea{Arithmetic sequences}
\LANGen{
\Question{
The \( 20 \)th term of an arithmetic sequence is \( 150 \) and the common difference is \( 3 \). Find the first term.}
{\( 93 \)}{\( 90 \)}{\( 87 \)}{\( 210 \)}{\( 207 \)}{}}
\LANGpl{
\Question{
\( 20 \) wyraz ciągu arytmetycznego to \( 150 \), a wspólna różnica wynosi \( 3 \). Oblicz pierwszy wyraz ciągu.}
{\( 93 \)}{\( 90 \)}{\( 87 \)}{\( 210 \)}{\( 207 \)}{}}
\LANGcs{
\Question{
Dvacátý člen aritmetické posloupnosti je roven \( 150 \) a její diference je \( 3 \). První člen se rovná:}
{\( 93 \)}{\( 90 \)}{\( 87 \)}{\( 210 \)}{\( 207 \)}{}}
\LANGsk{
\Question{
Dvadsiaty člen aritmetickej postupnosti je rovný \( 150 \) a jej diferencia je \( 3 \). Prvý člen sa rovná:}
{\( 93 \)}{\( 90 \)}{\( 87 \)}{\( 210 \)}{\( 207 \)}{}}
\LANGes{
\Question{
El término \( 20 \) de una progresión aritmética es \( 150 \) y la diferencia es \( 3 \). Halla el primer término.}
{\( 93 \)}{\( 90 \)}{\( 87 \)}{\( 210 \)}{\( 207 \)}{}}
\End

\Data\ID{1003057902}
\Updated{2020-07-15 03:07:12}
\Level{A}
\SubArea{Arithmetic sequences}
\LANGen{
\Question{
The \( 13 \)th term of an arithmetic sequence is \( 64 \) and the \( 21 \)st term is \( 32 \). Find the common difference.}
{\( -4 \)}{\( 4 \)}{\( 8 \)}{\( -8 \)}{\( 2 \)}{}}
\LANGpl{
\Question{
Trzynasty wyraz ciągu arytmetycznego to  \( 64 \), a dwudziesty pierwszy wyraz to \( 32 \). Znajdź wspólną różnice.}
{\( -4 \)}{\( 4 \)}{\( 8 \)}{\( -8 \)}{\( 2 \)}{}}
\LANGcs{
\Question{
Třináctý člen aritmetické posloupnosti je roven \( 64 \) a dvacátý první člen je \( 32 \). Diference této posloupnosti je:}
{\( -4 \)}{\( 4 \)}{\( 8 \)}{\( -8 \)}{\( 2 \)}{}}
\LANGsk{
\Question{
Trinásty člen aritmetickej postupnosti je rovný \( 64 \) a dvadsiaty prvý člen je \( 32 \). Diferencia tejto postupnosti je:}
{\( -4 \)}{\( 4 \)}{\( 8 \)}{\( -8 \)}{\( 2 \)}{}}
\LANGes{
\Question{
El término \( 13 \) de una progresión aritmética es \( 64 \) y el término \( 21 \) es \( 32 \). Halla la diferencia.}
{\( -4 \)}{\( 4 \)}{\( 8 \)}{\( -8 \)}{\( 2 \)}{}}
\End

\Data\ID{1003057903}
\Updated{2020-07-15 03:07:12}
\Level{A}
\SubArea{Arithmetic sequences}
\LANGen{
\Question{
The \( 17 \)th term of an arithmetic sequence is \( 77 \) and the common difference is \( 9 \). Choose the correct formula of calculation of the \( 5 \)th term.}
{\( a_5=77-12\cdot9 \)}{\(  a_5=17-12\cdot9 \)}{\( a_5=12\cdot9-77 \)}{\( a_5=77-16\cdot9 \)}{\( a_5=77+12\cdot9 \)}{}}
\LANGpl{
\Question{
\( 17 \) wyraz ciągu arytmetycznego to  \( 77 \) a wspólna różnica wynosi \( 9 \). Wybierz poprawny wzór na obliczenie  \( 5 \) wyrazu tego ciągu.}
{\( a_5=77-12\cdot9 \)}{\(  a_5=17-12\cdot9 \)}{\( a_5=12\cdot9-77 \)}{\( a_5=77-16\cdot9 \)}{\( a_5=77+12\cdot9 \)}{}}
\LANGcs{
\Question{
Sedmnáctý člen aritmetické posloupnosti je roven \( 77 \) a její diference je \( 9 \). Zvolte správný postup pro výpočet pátého členu.}
{\( a_5=77-12\cdot9 \)}{\(  a_5=17-12\cdot9 \)}{\( a_5=12\cdot9-77 \)}{\( a_5=77-16\cdot9 \)}{\( a_5=77+12\cdot9 \)}{}}
\LANGsk{
\Question{
Sedemnásty člen aritmetickej postupnosti je rovný \( 77 \) a jej diferencia je \( 9 \). Zvoľte správny postup pre výpočet piateho člena.}
{\( a_5=77-12\cdot9 \)}{\(  a_5=17-12\cdot9 \)}{\( a_5=12\cdot9-77 \)}{\( a_5=77-16\cdot9 \)}{\( a_5=77+12\cdot9 \)}{}}
\LANGes{
\Question{
El término \( 17 \) de una progresión aritmética es \( 77 \) y la diferencia es \( 9 \). Elige la fórmula correcta para hallar el quinto término.}
{\( a_5=77-12\cdot9 \)}{\(  a_5=17-12\cdot9 \)}{\( a_5=12\cdot9-77 \)}{\( a_5=77-16\cdot9 \)}{\( a_5=77+12\cdot9 \)}{}}
\End

\Data\ID{1003057904}
\Updated{2021-05-03 08:05:13}
\Level{A}
\SubArea{Arithmetic sequences}
\LANGen{
\Question{
The \( 2 \)nd term of an arithmetic sequence is \( 3 \) and the \( 14 \)th term is \( 51 \). Choose the incorrect formula:}
{\( a_{20}=a_2+(51-3)\cdot18 \)}{\( a_{20}=75 \)}{\( a_{20}=a_{14}+6\cdot4 \)}{\( a_{20}=a_2+18\cdot4 \)}{\( a_{20}=3+\frac{18}{12} (a_{14}-a_2 ) \)}{}}
\LANGpl{
\Question{
\( 2 \) wyraz ciągu arytmetycznego  to \( 3 \),a  \( 14 \) to \( 51 \). Wybierz wzór, który nie jest poprawny.}
{\( a_{20}=a_2+(51-3)\cdot18 \)}{\( a_{20}=75 \)}{\( a_{20}=a_{14}+6\cdot4 \)}{\( a_{20}=a_2+18\cdot4 \)}{\( a_{20}=3+\frac{18}{12} (a_{14}-a_2 ) \)}{}}
\LANGcs{
\Question{
Druhý člen aritmetické posloupnosti je roven \( 3 \) a čtrnáctý člen je \( 51 \). Vyberte tvrzení, které neplatí v této posloupnosti.}
{\( a_{20}=a_2+(51-3)\cdot18 \)}{\( a_{20}=75 \)}{\( a_{20}=a_{14}+6\cdot4 \)}{\( a_{20}=a_2+18\cdot4 \)}{\( a_{20}=3+\frac{18}{12} (a_{14}-a_2 ) \)}{}}
\LANGsk{
\Question{
Druhý člen aritmetickej postupnosti je rovný \( 3 \) a štrnásty člen je \( 51 \). Vyberte tvrdenie, ktoré neplatí v tejto postupnosti.}
{\( a_{20}=a_2+(51-3)\cdot18 \)}{\( a_{20}=75 \)}{\( a_{20}=a_{14}+6\cdot4 \)}{\( a_{20}=a_2+18\cdot4 \)}{\( a_{20}=3+\frac{18}{12} (a_{14}-a_2 ) \)}{}}
\LANGes{
\Question{
El segundo término de una progresión aritmética es \( 3 \) y el término \( 14 \) es \( 51 \). Elige la fórmula incorrecta:}
{\( a_{20}=a_2+(51-3)\cdot18 \)}{\( a_{20}=75 \)}{\( a_{20}=a_{14}+6\cdot4 \)}{\( a_{20}=a_2+18\cdot4 \)}{\( a_{20}=3+\frac{18}{12} (a_{14}-a_2 ) \)}{}}
\End

\Data\ID{1003057905}
\Updated{2020-07-15 03:07:12}
\Level{A}
\SubArea{Arithmetic sequences}
\LANGen{
\Question{
The \( n \)th term of an arithmetic sequence is \( 76 \), the common difference is \( 6 \) and the first term is \( -2 \). Find the \( n \).}
{\( 14 \)}{\( 12 \)}{\( 13 \)}{\( 15 \)}{\( 11 \)}{}}
\LANGpl{
\Question{
\( n \)-ty wyraz ciągu arytmetycznego to  \( 76 \), wspólna różnica wynosi \( 6 \), a pierwszy wyraz to \( -2 \). Znajdź \( n \).}
{\( 14 \)}{\( 12 \)}{\( 13 \)}{\( 15 \)}{\( 11 \)}{}}
\LANGcs{
\Question{
\( n \)-tý člen aritmetické posloupnosti je roven \( 76 \), diference je \( 6 \) a první člen je \( -2 \). Určete \( n \).}
{\( 14 \)}{\( 12 \)}{\( 13 \)}{\( 15 \)}{\( 11 \)}{}}
\LANGsk{
\Question{
V aritmetickej postupnosti je  \( n \)-tý člen rovný \( 76 \), diferencia je \( 6 \) a prvý člen je \( -2 \). Určte \( n \).}
{\( 14 \)}{\( 12 \)}{\( 13 \)}{\( 15 \)}{\( 11 \)}{}}
\LANGes{
\Question{
El término \( n \) de una progresión aritmética es \( 76 \), la diferencia es \( 6 \) y el primer término es \( -2 \). Halla el valor de \( n \).}
{\( 14 \)}{\( 12 \)}{\( 13 \)}{\( 15 \)}{\( 11 \)}{}}
\End

\Data\ID{1003085101}
\Updated{2021-05-03 08:05:46}
\Level{A}
\SubArea{Arithmetic sequences}
\LANGen{
\Question{
The second term of an arithmetic sequence is \( 3 \) and the fourth term is \( -1 \). Find the recursive formula for the sequence.}
{\( a_1=5;\ a_{n+1}=a_n-2  \text{ for all } n\in\mathbb{N} \)}{\( a_1=2;\ a_{n+1}=a_n-2  \text{ for all } n\in\mathbb{N} \)}{\( a_1=3;\ a_{n+1}=a_n-1  \text{ for all } n\in\mathbb{N} \)}{\( a_1=5;\ a_{n+1}=a_n-4  \text{ for all } n\in\mathbb{N} \)}{\( a_1=3;\ a_{n+1}=a_n-4  \text{ for all } n\in\mathbb{N} \)}{}}
\LANGpl{
\Question{
Drugim wyrazem ciągu arytmetycznego jest \( 3 \), a czwartym \( -1 \). Wyznacz wzór rekurencyjny tego ciągu.}
{\( a_1=5;\ a_{n+1}=a_n-2  \text{ dla wszystkich } n\in\mathbb{N} \)}{\( a_1=2;\ a_{n+1}=a_n-2  \text{ dla wszystkich } n\in\mathbb{N} \)}{\( a_1=3;\ a_{n+1}=a_n-1  \text{ dla wszystkich } n\in\mathbb{N} \)}{\( a_1=5;\ a_{n+1}=a_n-4  \text{ dla wszystkich } n\in\mathbb{N} \)}{\( a_1=3;\ a_{n+1}=a_n-4  \text{ dla wszystkich } n\in\mathbb{N} \)}{}}
\LANGcs{
\Question{
Najděte rekurentní vyjádření aritmetické posloupnosti, je-li druhý člen roven \( 3 \) a čtvrtý člen \( -1 \).}
{\( a_1=5;\ a_{n+1}=a_n-2,\ n\in\mathbb{N} \)}{\( a_1=2;\ a_{n+1}=a_n-2,\ n\in\mathbb{N} \)}{\( a_1=3;\ a_{n+1}=a_n-1,\ n\in\mathbb{N} \)}{\( a_1=5;\ a_{n+1}=a_n-4,\ n\in\mathbb{N} \)}{\( a_1=3;\ a_{n+1}=a_n-4,\ n\in\mathbb{N} \)}{}}
\LANGsk{
\Question{
Nájdite rekurentné vyjadrenie aritmetické postupnosti, ak jej druhý člen je rovný \( 3 \) a štvrtý člen \( -1 \).}
{\( a_1=5;\ a_{n+1}=a_n-2,\ n\in\mathbb{N} \)}{\( a_1=2;\ a_{n+1}=a_n-2,\ n\in\mathbb{N} \)}{\( a_1=3;\ a_{n+1}=a_n-1,\ n\in\mathbb{N} \)}{\( a_1=5;\ a_{n+1}=a_n-4,\ n\in\mathbb{N} \)}{\( a_1=3;\ a_{n+1}=a_n-4,\ n\in\mathbb{N} \)}{}}
\LANGes{
\Question{
El segundo término de una progresión aritmética es \( 3 \) y el cuarto término es \( -1 \). Halla la fórmula recursiva de la sucesión.}
{\( a_1=5;\ a_{n+1}=a_n-2  \text{ para todo } n\in\mathbb{N} \)}{\( a_1=2;\ a_{n+1}=a_n-2  \text{ para todo } n\in\mathbb{N} \)}{\( a_1=3;\ a_{n+1}=a_n-1  \text{ para todo } n\in\mathbb{N} \)}{\( a_1=5;\ a_{n+1}=a_n-4  \text{ para todo } n\in\mathbb{N} \)}{\( a_1=3;\ a_{n+1}=a_n-4  \text{ para todo } n\in\mathbb{N} \)}{}}
\End

\Data\ID{1003085102}
\Updated{2021-05-03 08:05:22}
\Level{A}
\SubArea{Arithmetic sequences}
\LANGen{
\Question{
The first term of an arithmetic sequence is \( 6 \) and the sixth  term is \( 1 \). Find the recursive formula for  the sequence.}
{\( a_1=6;\ a_{n+1}=a_n-1 \text{ for all } n\in\mathbb{N} \)}{\( a_1=6;\ a_{n+1}=a_n+1 \text{ for all } n\in\mathbb{N} \)}{\(a_1=1;\ a_{n+1}=a_n+5 \text{ for all } n\in\mathbb{N} \)}{\( a_1=1;\ a_{n+1}=a_n-5 \text{ for all } n\in\mathbb{N} \)}{}{}}
\LANGpl{
\Question{
Pierwszym wyrazem ciągu arytmetycznego jest \( 6 \), a szóstym \( 1 \). Wyznacz wzór rekurencyjny tego ciągu.}
{\( a_1=6;\ a_{n+1}=a_n-1 \text{ dla wszystkich } n\in\mathbb{N} \)}{\( a_1=6;\ a_{n+1}=a_n+1 \text{ dla wszystkich } n\in\mathbb{N} \)}{\(a_1=1;\ a_{n+1}=a_n+5 \text{ dla wszystkich } n\in\mathbb{N} \)}{\( a_1=1;\ a_{n+1}=a_n-5 \text{ dla wszystkich } n\in\mathbb{N} \)}{}{}}
\LANGcs{
\Question{
Najděte rekurentní vyjádření aritmetické posloupnosti, je-li první člen roven \( 6 \) a šestý člen \( 1 \).}
{\( a_1=6;\ a_{n+1}=a_n-1,\ n\in\mathbb{N} \)}{\( a_1=6;\ a_{n+1}=a_n+1,\ n\in\mathbb{N} \)}{\(a_1=1;\ a_{n+1}=a_n+5,\ n\in\mathbb{N} \)}{\( a_1=1;\ a_{n+1}=a_n-5,\ n\in\mathbb{N} \)}{}{}}
\LANGsk{
\Question{
Nájdite rekurentné vyjadrenie aritmetickej postupnosti, ak je prvý člen rovný \( 6 \) a šiesty člen \( 1 \).}
{\( a_1=6;\ a_{n+1}=a_n-1,\ n\in\mathbb{N} \)}{\( a_1=6;\ a_{n+1}=a_n+1,\ n\in\mathbb{N} \)}{\(a_1=1;\ a_{n+1}=a_n+5,\ n\in\mathbb{N} \)}{\( a_1=1;\ a_{n+1}=a_n-5,\ n\in\mathbb{N} \)}{}{}}
\LANGes{
\Question{
El primer término de una progresión aritmética es \( 6 \) y el sexto término es \( 1 \). Halla la fórmula recursiva de la sucesión.}
{\( a_1=6;\ a_{n+1}=a_n-1 \text{ para todo } n\in\mathbb{N} \)}{\( a_1=6;\ a_{n+1}=a_n+1 \text{ para todo } n\in\mathbb{N} \)}{\(a_1=1;\ a_{n+1}=a_n+5 \text{ para todo } n\in\mathbb{N} \)}{\( a_1=1;\ a_{n+1}=a_n-5 \text{ para todo } n\in\mathbb{N} \)}{}{}}
\End

\Data\ID{1003085103}
\Updated{2021-05-03 08:05:48}
\Level{A}
\SubArea{Arithmetic sequences}
\LANGen{
\Question{
The third term of an arithmetic sequence is \( 3 \) and the common difference is \( 3 \). Find the \(n\)th term.}
{\( a_n=3n-6 \text{ for all } n\in\mathbb{N} \)}{\( a_n=3n-3   \text{ for all } n\in\mathbb{N} \)}{\( a_n=3n  \text{ for all } n\in\mathbb{N} \)}{\( a_n=3n+3  \text{ for all } n\in\mathbb{N} \)}{\( a_n=3n+6   \text{ for all } n\in\mathbb{N} \)}{}}
\LANGpl{
\Question{
Trzecim wyrazem ciągu arytmetycznego jest \( 3 \) , a ogólna różnica to \( 3 \). Wyznacz \(n\)-ty wyraz tego ciągu.}
{\( a_n=3n-6 \text{ dla wszystkich } n\in\mathbb{N} \)}{\( a_n=3n-3   \text{ dla wszystkich } n\in\mathbb{N} \)}{\( a_n=3n  \text{ dla wszystkich } n\in\mathbb{N} \)}{\( a_n=3n+3  \text{ dla wszystkich} n\in\mathbb{N} \)}{\( a_n=3n+6   \text{ dla wszystkich  } n\in\mathbb{N} \)}{}}
\LANGcs{
\Question{
Najděte vzorec pro $n$-tý člen aritmetické posloupnosti, je-li třetí člen roven \( 3 \) a diference je \( 3 \).}
{\( a_n=3n-6;\ n\in\mathbb{N} \)}{\( a_n=3n-3;\ n\in\mathbb{N} \)}{\( a_n=3n;\ n\in\mathbb{N} \)}{\( a_n=3n+3;\ n\in\mathbb{N} \)}{\( a_n=3n+6;\ n\in\mathbb{N} \)}{}}
\LANGsk{
\Question{
Nájdite vzorec pre $n$-tý člen aritmetickej postupnosti, ak je tretí člen rovný \( 3 \) a diferencia je \( 3 \).}
{\( a_n=3n-6;\ n\in\mathbb{N} \)}{\( a_n=3n-3;\ n\in\mathbb{N} \)}{\( a_n=3n;\ n\in\mathbb{N} \)}{\( a_n=3n+3;\ n\in\mathbb{N} \)}{\( a_n=3n+6;\ n\in\mathbb{N} \)}{}}
\LANGes{
\Question{
El tercer término de una progresión aritmética es \( 3 \) y la diferencia es \( 3 \). Halla el término \(n\).}
{\( a_n=3n-6 \text{ para todo } n\in\mathbb{N} \)}{\( a_n=3n-3   \text{ para todo } n\in\mathbb{N} \)}{\( a_n=3n  \text{ para todo } n\in\mathbb{N} \)}{\( a_n=3n+3  \text{ para todo } n\in\mathbb{N} \)}{\( a_n=3n+6   \text{ para todo } n\in\mathbb{N} \)}{}}
\End

\Data\ID{1003085104}
\Updated{2020-07-15 03:07:21}
\Level{A}
\SubArea{Arithmetic sequences}
\LANGen{
\Question{
The \( n \)th term of an arithmetic sequence is \( 1-3n \). Find the \( 5 \)th term and the common difference.}
{\( a_5=-14;\ d=-3 \)}{\( a_5=-2;\ d=-3 \)}{\( a_5=14;\ d=-3 \)}{\( a_5=-14;\ d=3 \)}{\( a_5=-2;\ d=3 \)}{}}
\LANGpl{
\Question{
\( n \)-tym wyrazem ciągu arytmetycznego jest \( 1-3n \). Wyznacz \( 5 \)-ty wyraz i ogólną różnicę. .}
{\( a_5=-14;\ d=-3 \)}{\( a_5=-2;\ d=-3 \)}{\( a_5=14;\ d=-3 \)}{\( a_5=-14;\ d=3 \)}{\( a_5=-2;\ d=3 \)}{}}
\LANGcs{
\Question{
Určete pátý člen a diferenci aritmetické posloupnosti, je-li $n$-tý člen roven \( 1-3n;\ n\in\mathbb{N} \).}
{\( a_5=-14;\ d=-3 \)}{\( a_5=-2;\ d=-3 \)}{\( a_5=14;\ d=-3 \)}{\( a_5=-14;\ d=3 \)}{\( a_5=-2;\ d=3 \)}{}}
\LANGsk{
\Question{
Určte piaty člen a diferenciu aritmetickej postupnosti, ak je $n$-tý člen rovný \( 1-3n;\ n\in\mathbb{N} \).}
{\( a_5=-14;\ d=-3 \)}{\( a_5=-2;\ d=-3 \)}{\( a_5=14;\ d=-3 \)}{\( a_5=-14;\ d=3 \)}{\( a_5=-2;\ d=3 \)}{}}
\LANGes{
\Question{
El término \( n \) de una progresión aritmética es \( 1-3n \). Halla el quinto término y la diferencia.}
{\( a_5=-14;\ d=-3 \)}{\( a_5=-2;\ d=-3 \)}{\( a_5=14;\ d=-3 \)}{\( a_5=-14;\ d=3 \)}{\( a_5=-2;\ d=3 \)}{}}
\End

\Data\ID{1003085208}
\Updated{2020-12-28 01:12:45}
\Level{A}
\SubArea{Arithmetic sequences}
\LANGen{
\Question{
Given an arithmetic sequence $(a_n)_{n=1}^{\infty}$, where $a_8=\frac83$, $a_k=8$ and the common difference is $\frac23$, find $k$.}
{$16$}{$15$}{$14$}{$21$}{$10$}{}}
\LANGpl{
\Question{
Dany jest ciąg arytmetyczny $\{a_n\}_{n=1}^{\infty}$, gdzie $a_8=\frac83$, $a_k=8$, a powszechna różnica wynosi $\frac23$, wyznacz $k$.}
{$16$}{$15$}{$14$}{$21$}{$10$}{}}
\LANGcs{
\Question{
V aritmetické posloupnosti je dáno: $a_8=\frac83$, $a_k=8$ a diference je $\frac23$. Určete $k$.}
{$16$}{$15$}{$14$}{$21$}{$10$}{}}
\LANGsk{
\Question{
V aritmetickej postupnosti $(a_n)_{n=1}^{\infty}$ je dané: $a_8=\frac83$, $a_k=8$ a diferencia je $\frac23$. Určte $k$.}
{$16$}{$15$}{$14$}{$21$}{$10$}{}}
\LANGes{
\Question{
Dada la progresión aritmética $\{a_n\}_{n=1}^{\infty}$, donde $a_8=\frac83$, $a_k=8$ y la diferencia es $\frac23$, halla $k$.}
{$16$}{$15$}{$14$}{$21$}{$10$}{}}
\End

\Data\ID{9000065301}
\Updated{2020-07-15 03:07:37}
\Level{A}
\SubArea{Arithmetic sequences}
\LANGen{
\Question{
Find the recurrence equations for the arithmetic sequence with the first term
\(a_{1} = 4\) and the common
difference \(d = -2\).}
{\(a_{1} = 4;\ a_{n+1} = a_{n} - 2,\ n\in\mathbb{N}\)}{\(a_{1} = 4;\ a_{n+1} = a_{1} - 2,\ n\in\mathbb{N}\)}{\(a_{n} = 4 + a_{n+2},\ n\in\mathbb{N}\)}{\(a_{n+1} = a_{n} + 2,\ n\in\mathbb{N}\)}{}{}}
\LANGpl{
\Question{
Wyznacz równania rekurencyjne dla ciągu arytmetycznego, w którym pierwszy wyraz jest równy 
\(a_{1} = 4\), a różnica \(d = -2\).}
{\(a_{1} = 4;\ a_{n+1} = a_{n} - 2,\ n\in\mathbb{N}\)}{\(a_{1} = 4;\ a_{n+1} = a_{1} - 2,\ n\in\mathbb{N}\)}{\(a_{n} = 4 + a_{n+2},\ n\in\mathbb{N}\)}{\(a_{n+1} = a_{n} + 2,\ n\in\mathbb{N}\)}{}{}}
\LANGcs{
\Question{
Najděte rekurentní vyjádření aritmetické posloupnosti, je-li dáno
\(a_{1} = 4\),
\(d = -2\).}
{\(a_{1} = 4;\ a_{n+1} = a_{n} - 2,\ n\in\mathbb{N}\)}{\(a_{1} = 4;\ a_{n+1} = a_{1} - 2,\ n\in\mathbb{N}\)}{\(a_{n} = 4 + a_{n+2},\ n\in\mathbb{N}\)}{\(a_{n+1} = a_{n} + 2,\ n\in\mathbb{N}\)}{}{}}
\LANGsk{
\Question{
Nájdite rekurentné vyjadrenie aritmetickej postupnosti, ak je dané
\(a_{1} = 4\),
\(d = -2\).}
{\(a_{1} = 4;\ a_{n+1} = a_{n} - 2,\ n\in\mathbb{N}\)}{\(a_{1} = 4;\ a_{n+1} = a_{1} - 2,\ n\in\mathbb{N}\)}{\(a_{n} = 4 + a_{n+2},\ n\in\mathbb{N}\)}{\(a_{n+1} = a_{n} + 2,\ n\in\mathbb{N}\)}{}{}}
\LANGes{
\Question{
Suponiendo que \(a_{1} = 4\) y \(d = -2\), halla la fórmula recursiva de la progresión aritmética.}
{\(a_{1} = 4;\ a_{n+1} = a_{n} - 2,\ n\in\mathbb{N}\)}{\(a_{1} = 4;\ a_{n+1} = a_{1} - 2,\ n\in\mathbb{N}\)}{\(a_{n} = 4 + a_{n+2},\ n\in\mathbb{N}\)}{\(a_{n+1} = a_{n} + 2,\ n\in\mathbb{N}\)}{}{}}
\End

\Data\ID{9000065302}
\Updated{2020-07-15 03:07:37}
\Level{A}
\SubArea{Arithmetic sequences}
\LANGen{
\Question{
Find the formula for the \(n\)-th
term of an arithmetic sequence with the first term
\(a_{1} = 1\) and the
second term \(a_{2} = -2\).}
{\(a_{n} = 4 - 3n,\ n\in\mathbb{N}\)}{\(a_{n} = 1 - 2n,\ n\in\mathbb{N}\)}{\(a_{n} = -2 + n,\ n\in\mathbb{N}\)}{\(a_{n} = 3 + 2n,\ n\in\mathbb{N}\)}{}{}}
\LANGpl{
\Question{
Wyznacz wzór dla \(n\)-tego 
wyrazu ciągu arytmetycznego, gdy pierwszy wyraz tego ciągu wynosi 
\(a_{1} = 1\), a drugi jest równy  \(a_{2} = -2\).}
{\(a_{n} = 4 - 3n,\ n\in\mathbb{N}\)}{\(a_{n} = 1 - 2n,\ n\in\mathbb{N}\)}{\(a_{n} = -2 + n,\ n\in\mathbb{N}\)}{\(a_{n} = 3 + 2n,\ n\in\mathbb{N}\)}{}{}}
\LANGcs{
\Question{
Najděte vzorec pro \(n\)-tý
člen aritmetické posloupnosti, je-li dáno
\(a_{1} = 1\),
\(a_{2} = -2\).}
{\(a_{n} = 4 - 3n,\ n\in\mathbb{N}\)}{\(a_{n} = 1 - 2n,\ n\in\mathbb{N}\)}{\(a_{n} = -2 + n,\ n\in\mathbb{N}\)}{\(a_{n} = 3 + 2n,\ n\in\mathbb{N}\)}{}{}}
\LANGsk{
\Question{
Nájdite vzorec pre \(n\)-tý
člen aritmetickej postupnosti, ak je dané
\(a_{1} = 1\),
\(a_{2} = -2\).}
{\(a_{n} = 4 - 3n,\ n\in\mathbb{N}\)}{\(a_{n} = 1 - 2n,\ n\in\mathbb{N}\)}{\(a_{n} = -2 + n,\ n\in\mathbb{N}\)}{\(a_{n} = 3 + 2n,\ n\in\mathbb{N}\)}{}{}}
\LANGes{
\Question{
Halla el término \(n\) de una progresión aritmética con el primer término
\(a_{1} = 1\) y el segundo término \(a_{2} = -2\).}
{\(a_{n} = 4 - 3n,\ n\in\mathbb{N}\)}{\(a_{n} = 1 - 2n,\ n\in\mathbb{N}\)}{\(a_{n} = -2 + n,\ n\in\mathbb{N}\)}{\(a_{n} = 3 + 2n,\ n\in\mathbb{N}\)}{}{}}
\End

\Data\ID{9000065303}
\Updated{2020-07-15 03:07:37}
\Level{A}
\SubArea{Arithmetic sequences}
\LANGen{
\Question{
Find the recurrence equations for the arithmetic sequence with the second term
\(a_{2} = 7\) and the common
difference \(d = 4\).}
{\(a_{1} = 3;\ a_{n} = a_{n-1} + 4,\ n\in\mathbb{N}\)}{\(a_{1} = 7;\ a_{n+1} = a_{n} + 4,\ n\in\mathbb{N}\)}{\(a_{n} = 7 + a_{n+4},\ n\in\mathbb{N}\)}{\(a_{n+1} = a_{n} + 7,\ n\in\mathbb{N}\)}{}{}}
\LANGpl{
\Question{
Wyznacz równania rekurencyjne dla ciągu arytmetycznego, w którym drugi wyraz jest równy 
\(a_{2} = 7\), a różnica \(d = 4\).}
{\(a_{1} = 3;\ a_{n} = a_{n-1} + 4,\ n\in\mathbb{N}\)}{\(a_{1} = 7;\ a_{n+1} = a_{n} + 4,\ n\in\mathbb{N}\)}{\(a_{n} = 7 + a_{n+4},\ n\in\mathbb{N}\)}{\(a_{n+1} = a_{n} + 7,\ n\in\mathbb{N}\)}{}{}}
\LANGcs{
\Question{
Najděte rekurentní vyjádření aritmetické posloupnosti, je-li dáno
\(a_{2} = 7\),
\(d = 4\).}
{\(a_{1} = 3;\ a_{n} = a_{n-1} + 4,\ n\in\mathbb{N}\)}{\(a_{1} = 7;\ a_{n+1} = a_{n} + 4,\ n\in\mathbb{N}\)}{\(a_{n} = 7 + a_{n+4},\ n\in\mathbb{N}\)}{\(a_{n+1} = a_{n} + 7,\ n\in\mathbb{N}\)}{}{}}
\LANGsk{
\Question{
Nájdite rekurentné vyjadrenie aritmetickej postupnosti, ak je dané
\(a_{2} = 7\),
\(d = 4\).}
{\(a_{1} = 3;\ a_{n} = a_{n-1} + 4,\ n\in\mathbb{N}\)}{\(a_{1} = 7;\ a_{n+1} = a_{n} + 4,\ n\in\mathbb{N}\)}{\(a_{n} = 7 + a_{n+4},\ n\in\mathbb{N}\)}{\(a_{n+1} = a_{n} + 7,\ n\in\mathbb{N}\)}{}{}}
\LANGes{
\Question{
Suponiendo que
\(a_{2} = 7\) y \(d = 4\), halla la fórmula recursiva de la progresión aritmética.}
{\(a_{1} = 3;\ a_{n} = a_{n-1} + 4,\ n\in\mathbb{N}\)}{\(a_{1} = 7;\ a_{n+1} = a_{n} + 4,\ n\in\mathbb{N}\)}{\(a_{n} = 7 + a_{n+4},\ n\in\mathbb{N}\)}{\(a_{n+1} = a_{n} + 7,\ n\in\mathbb{N}\)}{}{}}
\End

\Data\ID{9000065304}
\Updated{2020-07-15 03:07:37}
\Level{A}
\SubArea{Arithmetic sequences}
\LANGen{
\Question{
Find the first term \(a_{1}\) and
the common difference \(d\) of the
arithmetic sequence \((5 + 2n)_{n=1}^{\infty }\).}
{\(a_{1} = 7;\ d = 2\)}{\(a_{1} = 5;\ d = 2\)}{\(a_{1} = 3;\ d = -2\)}{\(a_{1} = 2;\ d = 5\)}{}{}}
\LANGpl{
\Question{
Jaki jest pierwszy wyraz \(a_{1}\) i
i różnica  \(d\) ciągu arytmetycznego \((5 + 2n)_{n=1}^{\infty }\)?}
{\(a_{1} = 7;\ d = 2\)}{\(a_{1} = 5;\ d = 2\)}{\(a_{1} = 3;\ d = -2\)}{\(a_{1} = 2;\ d = 5\)}{}{}}
\LANGcs{
\Question{
Určete první člen a diferenci aritmetické posloupnosti
\((5 + 2n)_{n=1}^{\infty }\).}
{\(a_{1} = 7;\ d = 2\)}{\(a_{1} = 5;\ d = 2\)}{\(a_{1} = 3;\ d = -2\)}{\(a_{1} = 2;\ d = 5\)}{}{}}
\LANGsk{
\Question{
Určte prvý člen a diferenciu aritmetickej postupnosti
\((5 + 2n)_{n=1}^{\infty }\).}
{\(a_{1} = 7;\ d = 2\)}{\(a_{1} = 5;\ d = 2\)}{\(a_{1} = 3;\ d = -2\)}{\(a_{1} = 2;\ d = 5\)}{}{}}
\LANGes{
\Question{
Halla el primer término \(a_{1}\) y la diferencia \(d\) de la progresión aritmética \((5 + 2n)_{n=1}^{\infty }\).}
{\(a_{1} = 7;\ d = 2\)}{\(a_{1} = 5;\ d = 2\)}{\(a_{1} = 3;\ d = -2\)}{\(a_{1} = 2;\ d = 5\)}{}{}}
\End

\Data\ID{9000065305}
\Updated{2021-05-03 08:05:58}
\Level{A}
\SubArea{Arithmetic sequences}
\LANGen{
\Question{
In the arithmetic sequence given by the relations
\(a_{1} =\pi \),
\(a_{n+1} = a_{n} + 2\pi \) find
\(a_{13}\).}
{\(a_{13} = 25\pi \)}{\(a_{13} = 27\pi \)}{\(a_{13} = 26\pi \)}{\(a_{13} = 24\pi \)}{}{}}
\LANGpl{
\Question{
Dany jest ciąg arytmetyczny, gdzie
\(a_{1} =\pi \),
\(a_{n+1} = a_{n} + 2\pi \). Określ 
\(a_{13}\).}
{\(a_{13} = 25\pi \)}{\(a_{13} = 27\pi \)}{\(a_{13} = 26\pi \)}{\(a_{13} = 24\pi \)}{}{}}
\LANGcs{
\Question{
Určete třináctý člen aritmetické posloupnosti, je-li dáno
\(a_{1} =\pi \),
\(a_{n+1} = a_{n} + 2\pi \).}
{\(a_{13} = 25\pi \)}{\(a_{13} = 27\pi \)}{\(a_{13} = 26\pi \)}{\(a_{13} = 24\pi \)}{}{}}
\LANGsk{
\Question{
Určte trinásty člen aritmetickej postupnosti, ak je dané
\(a_{1} =\pi \),
\(a_{n+1} = a_{n} + 2\pi \).}
{\(a_{13} = 25\pi \)}{\(a_{13} = 27\pi \)}{\(a_{13} = 26\pi \)}{\(a_{13} = 24\pi \)}{}{}}
\LANGes{
\Question{
Suponiendo que
\(a_{1} =\pi \) y
\(a_{n+1} = a_{n} + 2\pi \), y sabiendo que se trata de una progresión aritmética, halla
\(a_{13}\).}
{\(a_{13} = 25\pi \)}{\(a_{13} = 27\pi \)}{\(a_{13} = 26\pi \)}{\(a_{13} = 24\pi \)}{}{}}
\End

\Data\ID{9000065306}
\Updated{2021-05-03 08:05:53}
\Level{A}
\SubArea{Arithmetic sequences}
\LANGen{
\Question{
In the arithmetic sequence given by the second term
\(a_{2} = -3\) and the
fifth term \(a_{5} = 3\) 
find \(a_{11}\).}
{\(a_{11} = 15\)}{\(a_{11} = 22\)}{\(a_{11} = 19\)}{\(a_{11} = 27\)}{}{}}
\LANGpl{
\Question{
Dany jest ciąg arytmetyczny, gdzie drugi wyraz jest równy
\(a_{2} = -3\), a piąty  \(a_{5} = 3\). Oblicz  \(a_{11}\).}
{\(a_{11} = 15\)}{\(a_{11} = 22\)}{\(a_{11} = 19\)}{\(a_{11} = 27\)}{}{}}
\LANGcs{
\Question{
Určete jedenáctý člen aritmetické posloupnosti, je-li dáno
\(a_{2} = -3\),
\(a_{5} = 3\).}
{\(a_{11} = 15\)}{\(a_{11} = 22\)}{\(a_{11} = 19\)}{\(a_{11} = 27\)}{}{}}
\LANGsk{
\Question{
Určte jedenásty člen aritmetickej postupnosti, ak je dané
\(a_{2} = -3\),
\(a_{5} = 3\).}
{\(a_{11} = 15\)}{\(a_{11} = 22\)}{\(a_{11} = 19\)}{\(a_{11} = 27\)}{}{}}
\LANGes{
\Question{
Suponiendo que
\(a_{2} = -3\) y \(a_{5} = 3\), y sabiendo que se trata de una progresión aritmética,
halla \(a_{11}\).}
{\(a_{11} = 15\)}{\(a_{11} = 22\)}{\(a_{11} = 19\)}{\(a_{11} = 27\)}{}{}}
\End

\Data\ID{9000065309}
\Updated{2021-05-03 08:05:08}
\Level{A}
\SubArea{Arithmetic sequences}
\LANGen{
\Question{
The arithmetic sequence satisfies \(a_{26} = 58\) 
and \(a_{21} = 43\). Find the
first term \(a_{1}\) and
the common difference \(d\) 
of this sequence.}
{\(a_{1} = -17;\ d = 3\)}{\(a_{1} = -1;\ d = 5\)}{\(a_{1} = 1;\ d = 15\)}{\(a_{1} = -1;\ d = 3\)}{}{}}
\LANGpl{
\Question{
Dany jest ciąg arytmetyczny, gdzie  \(a_{26} = 58\), a \(a_{21} = 43\). Oblicz pierwszy wyraz tego ciągu  \(a_{1}\) i
różnicę \(d\).}
{\(a_{1} = -17;\ d = 3\)}{\(a_{1} = -1;\ d = 5\)}{\(a_{1} = 1;\ d = 15\)}{\(a_{1} = -1;\ d = 3\)}{}{}}
\LANGcs{
\Question{
Určete první člen a diferenci aritmetické posloupnosti, je-li dáno
\(a_{26} = 58\),
\(a_{21} = 43\).}
{\(a_{1} = -17;\ d = 3\)}{\(a_{1} = -1;\ d = 5\)}{\(a_{1} = 1;\ d = 15\)}{\(a_{1} = -1;\ d = 3\)}{}{}}
\LANGsk{
\Question{
Určte prvý člen a diferenciu aritmetickej postupnosti, ak je dané
\(a_{26} = 58\),
\(a_{21} = 43\).}
{\(a_{1} = -17;\ d = 3\)}{\(a_{1} = -1;\ d = 5\)}{\(a_{1} = 1;\ d = 15\)}{\(a_{1} = -1;\ d = 3\)}{}{}}
\LANGes{
\Question{
Suponiendo que \(a_{26} = 58\) 
y \(a_{21} = 43\), halla el primer término \(a_{1}\) y la diferencia \(d\) sabiendo que se trata de una sucesión aritmética.}
{\(a_{1} = -17;\ d = 3\)}{\(a_{1} = -1;\ d = 5\)}{\(a_{1} = 1;\ d = 15\)}{\(a_{1} = -1;\ d = 3\)}{}{}}
\End

\Data\ID{1003057906}
\Updated{2020-07-15 03:07:12}
\Level{B}
\SubArea{Arithmetic sequences}
\LANGen{
\Question{
The sum of the first \( 15 \) terms of an arithmetic sequence is \( 210 \) and the \( 15 \)th term is \( 7 \). Choose the correct formula of calculation of the first term.}
{\( a_1=\frac2{15}\cdot210-7 \)}{\(  a_1=\frac2{15}\cdot\frac{210}7 \)}{\( a_1=\frac{15}2\cdot210-7 \)}{\( a_1= \frac{210}7 \)}{\( a_1=\frac2{15} (210-7) \)}{}}
\LANGpl{
\Question{
Suma pierwszych \( 15 \) wyrazów ciągu arytmetycznego wynosi \( 210 \), a  \( 15 \)-ty wyraz ciągu to \( 7 \). Wybierz poprawny wzór na obliczenie wartości pierwszego wyrazu tego ciągu.}
{\( a_1=\frac2{15}\cdot210-7 \)}{\(  a_1=\frac2{15}\cdot\frac{210}7 \)}{\( a_1=\frac{15}2\cdot210-7 \)}{\( a_1= \frac{210}7 \)}{\( a_1=\frac2{15} (210-7) \)}{}}
\LANGcs{
\Question{
Součet prvních patnácti členů aritmetické posloupnosti je \( 210 \) a patnáctý člen je roven \( 7 \). Pro první člen této posloupnosti platí:}
{\( a_1=\frac2{15}\cdot210-7 \)}{\(  a_1=\frac2{15}\cdot\frac{210}7 \)}{\( a_1=\frac{15}2\cdot210-7 \)}{\( a_1= \frac{210}7 \)}{\( a_1=\frac2{15} (210-7) \)}{}}
\LANGsk{
\Question{
Súčet prvých pätnástich členov aritmetickej postupnosti je \( 210 \) a pätnásty člen je rovný \( 7 \). Pre prvý člen tejto postupnosti platí:}
{\( a_1=\frac2{15}\cdot210-7 \)}{\(  a_1=\frac2{15}\cdot\frac{210}7 \)}{\( a_1=\frac{15}2\cdot210-7 \)}{\( a_1= \frac{210}7 \)}{\( a_1=\frac2{15} (210-7) \)}{}}
\LANGes{
\Question{
La suma de los \( 15 \) primeros términos de una progresión aritmética es \( 210 \) y el término \( 15 \) es \( 7 \). Elige la fórmula correcta para hallar el primer término.}
{\( a_1=\frac2{15}\cdot210-7 \)}{\(  a_1=\frac2{15}\cdot\frac{210}7 \)}{\( a_1=\frac{15}2\cdot210-7 \)}{\( a_1= \frac{210}7 \)}{\( a_1=\frac2{15} (210-7) \)}{}}
\End

\Data\ID{1003057907}
\Updated{2020-12-28 02:12:32}
\Level{B}
\SubArea{Arithmetic sequences}
\LANGen{
\Question{
The sum of the first \( 25 \) terms of an arithmetic sequence is \( 700 \) and the first term is \( 4 \). Choose the incorrect statement about its common difference.}
{\( d \) is an odd number}{\( d < 4 \)}{\( d > 0 \)}{\( d \) is a divisor of \( 48 \)}{}{}}
\LANGpl{
\Question{
Suma pierwszych \( 25 \) wyrazów ciągu arytmetycznego wynosi  \( 700 \), a pierwszy wyraz to \( 4 \). Wybierz niepoprawny wzór na oszacowanie ich wspólnej różnicy.}
{\( d \) jest liczbą nieparzystą}{\( d < 4 \)}{\( d > 0 \)}{\( d \) jest dzielnikiem \( 48 \)}{}{}}
\LANGcs{
\Question{
Součet prvních \( 25 \) členů aritmetické posloupnosti je \( 700 \) a první člen je \( 4 \). Pro diferenci této posloupnosti neplatí:}
{\( d \) je liché číslo}{\( d < 4 \)}{\( d > 0 \)}{\( d \) je dělitel čísla \( 48 \)}{}{}}
\LANGsk{
\Question{
Súčet prvých \( 25 \) členov aritmetickej postupnosti je \( 700 \) a prvý člen je \( 4 \). Vyberte nepravdivé tvrdenie o diferencii tejto postupnosti.}
{\( d \) je nepárne číslo}{\( d < 4 \)}{\( d > 0 \)}{\( d \) je deliteľ čísla \( 48 \)}{}{}}
\LANGes{
\Question{
La suma de los \( 25 \) primeros términos de una progresión aritmética es \( 700 \) y el primer término es \( 4 \). Elige la fórmula incorrecta para hallar su diferencia.}
{\( d \) es un número impar}{\( d < 4 \)}{\( d > 0 \)}{\( d \) es un divisor de \( 48 \)}{}{}}
\End

\Data\ID{1003057908}
\Updated{2020-07-15 03:07:12}
\Level{B}
\SubArea{Arithmetic sequences}
\LANGen{
\Question{
The \( 3 \)rd term of an arithmetic sequence is \( 2 \) and the \( 20 \)th term is \( 53 \). Find the sum of the first \( 13 \) terms.}
{\( 182 \)}{\( 358 \)}{\( 364 \)}{\( 689 \)}{\( 106 \)}{}}
\LANGpl{
\Question{
\( 3 \) wyraz ciągu arytmetycznego to  \( 2 \), a  \( 20 \)ty to \( 53 \). Oblicz sumę pierwszych  \( 13 \) wyrazów tego ciągu.}
{\( 182 \)}{\( 358 \)}{\( 364 \)}{\( 689 \)}{\( 106 \)}{}}
\LANGcs{
\Question{
Třetí člen aritmetické posloupnosti je \( 2 \) a dvacátý člen je \( 53 \). Součet prvních třinácti členů této posloupnosti je:}
{\( 182 \)}{\( 358 \)}{\( 364 \)}{\( 689 \)}{\( 106 \)}{}}
\LANGsk{
\Question{
Tretí člen aritmetickej postupnosti je \( 2 \) a dvadsiaty člen je \( 53 \). Súčet prvých trinásť členov tejto postupnosti je:}
{\( 182 \)}{\( 358 \)}{\( 364 \)}{\( 689 \)}{\( 106 \)}{}}
\LANGes{
\Question{
El tercer término de una progresión aritmética es \( 2 \) y el término \( 20 \) es \( 53 \). Halla la suma de los \( 13 \) primeros términos.}
{\( 182 \)}{\( 358 \)}{\( 364 \)}{\( 689 \)}{\( 106 \)}{}}
\End

\Data\ID{1003057909}
\Updated{2020-11-12 04:11:48}
\Level{B}
\SubArea{Arithmetic sequences}
\LANGen{
\Question{
The \( 2 \)nd term of an arithmetic sequence is \( 100 \) and the common difference is \( -2 \). Choose the correct statement about the sum of the first \( 100 \) terms.}
{\( s_{100} > 200 \)}{\( s_{100} > 0 \) and \( s_{100} < 200 \)}{\( s_{100} < 0 \) and \( s_{100} > -100 \)}{\( s_{100} < -100 \)}{\( s_{100}=0 \)}{}}
\LANGpl{
\Question{
\( 2 \) wyraz ciągu arytmetycznego to \( 100 \), a wspólna różnica to \( -2 \). Wybierz poprawny wzór na oszacowanie sumy pierwszych \( 100 \) wyrazów tego ciągu.}
{\( s_{100} > 200 \)}{\( s_{100} > 0 \) and \( s_{100} < 200 \)}{\( s_{100} < 0 \) and \( s_{100} > -100 \)}{\( s_{100} < -100 \)}{\( s_{100}=0 \)}{}}
\LANGcs{
\Question{
Druhý člen aritmetické posloupnosti je \( 100 \) a její diference je \( -2 \). Pro součet prvních \( 100 \) členů této posloupnosti platí:}
{\( s_{100} > 200 \)}{\( s_{100} > 0 \) a \( s_{100} < 200 \)}{\( s_{100} < 0 \) a \( s_{100} > -100 \)}{\( s_{100} < -100 \)}{\( s_{100}=0 \)}{}}
\LANGsk{
\Question{
Druhý člen aritmetickej postupnosti je \( 100 \) a jej diferencia je \( -2 \). Pre súčet prvých \( 100 \) členov tejto postupnosti platí:}
{\( s_{100} > 200 \)}{\( s_{100} > 0 \) a \( s_{100} < 200 \)}{\( s_{100} < 0 \) a \( s_{100} > -100 \)}{\( s_{100} < -100 \)}{\( s_{100}=0 \)}{}}
\LANGes{
\Question{
El segundo término de una progresión aritmética es \( 100 \) y la diferencia es \( -2 \). Elige la fórmula correcta para hallar la suma de los \( 100 \) primeros términos.}
{\( s_{100} > 200 \)}{\( s_{100} > 0 \) and \( s_{100} < 200 \)}{\( s_{100} < 0 \) and \( s_{100} > -100 \)}{\( s_{100} < -100 \)}{\( s_{100}=0 \)}{}}
\End

\Data\ID{1003057910}
\Updated{2020-07-15 03:07:12}
\Level{B}
\SubArea{Arithmetic sequences}
\LANGen{
\Question{
The sum of the first \( n \) terms of an arithmetic sequence is \( 0 \), the common difference is \( 3 \) and the first term is \( -45 \). Find the \( n \).}
{\( 31 \)}{\( 15 \)}{\( 30 \)}{\( 16 \)}{\( 32 \)}{}}
\LANGpl{
\Question{
Suma pierwszych \( n \) wyrazów ciągu arytmetycznego wynosi \( 0 \), wspólna różnica \( 3 \), a pierwszy wyraz to \( -45 \). Znajdź wartość \( n \).}
{\( 31 \)}{\( 15 \)}{\( 30 \)}{\( 16 \)}{\( 32 \)}{}}
\LANGcs{
\Question{
Součet prvních \( n \) členů aritmetické posloupnosti je \( 0 \), diference je rovna \( 3 \) a první člen je \( -45 \). Určete \( n \).}
{\( 31 \)}{\( 15 \)}{\( 30 \)}{\( 16 \)}{\( 32 \)}{}}
\LANGsk{
\Question{
Súčet prvých \( n \) členov aritmetickej postupnosti je \( 0 \), diferencia je rovná \( 3 \) a prvý člen je \( -45 \). Určte \( n \).}
{\( 31 \)}{\( 15 \)}{\( 30 \)}{\( 16 \)}{\( 32 \)}{}}
\LANGes{
\Question{
La suma de los \( n \) primeros términos de una progresión aritmética es \( 0 \), la diferencia es \( 3 \) y el primer término es \( -45 \). Halla el valor de \( n \).}
{\( 31 \)}{\( 15 \)}{\( 30 \)}{\( 16 \)}{\( 32 \)}{}}
\End

\Data\ID{1003085201}
\Updated{2020-07-15 03:07:12}
\Level{B}
\SubArea{Arithmetic sequences}
\LANGen{
\Question{
Find the sum of all even numbers which are greater than $6$ and smaller than $122$.}
{$3648$}{$3584$}{$3654$}{$3712$}{$3776$}{}}
\LANGpl{
\Question{
Wyznacz sumę wszystkich parzystych liczb większych niż $6$ i mniejszych niż $122$.}
{$3648$}{$3584$}{$3654$}{$3712$}{$3776$}{}}
\LANGcs{
\Question{
Určete součet všech sudých čísel větších než $6$ a menších než $122$.}
{$3648$}{$3584$}{$3654$}{$3712$}{$3776$}{}}
\LANGsk{
\Question{
Určte súčet všetkých párnych čísel väčších než $6$ a menších než $122$.}
{$3648$}{$3584$}{$3654$}{$3712$}{$3776$}{}}
\LANGes{
\Question{
Halla la suma de todos los números pares mayores que $6$ y menores que $122$.}
{$3648$}{$3584$}{$3654$}{$3712$}{$3776$}{}}
\End

\Data\ID{1003085202}
\Updated{2020-10-02 10:10:10}
\Level{B}
\SubArea{Arithmetic sequences}
\LANGen{
\Question{
What is the sum of all two-digit numbers divisible by five?}
{$945$}{$950$}{$1050$}{$1045$}{$940$}{}}
\LANGpl{
\Question{
Jaka jest suma wszystkich dwucyfrowych liczb podzielnych przez pięć?}
{$945$}{$950$}{$1050$}{$1045$}{$940$}{}}
\LANGcs{
\Question{
Součet všech dvojciferných čísel dělitelných pěti je:}
{$945$}{$950$}{$1050$}{$1045$}{$940$}{}}
\LANGsk{
\Question{
Súčet všetkých dvojciferných čísel deliteľných piatimi je:}
{$945$}{$950$}{$1050$}{$1045$}{$940$}{}}
\LANGes{
\Question{
Halla la suma de todos los números de dos dígitos y divisibles por cinco.}
{$945$}{$950$}{$1050$}{$1045$}{$940$}{}}
\End

\Data\ID{1003085203}
\Updated{2020-10-02 10:10:47}
\Level{B}
\SubArea{Arithmetic sequences}
\LANGen{
\Question{
What is the sum of all positive multiples of $7$ smaller than $100$?}
{$735$}{$637$}{$728$}{$105$}{$686$}{}}
\LANGpl{
\Question{
Jaka jest suma wszystkich dodatnich wielokrotności  $7$ mniejszych niż $100$?}
{$735$}{$637$}{$728$}{$105$}{$686$}{}}
\LANGcs{
\Question{
Součet všech kladných násobků čísla $7$ menších než $100$ je:}
{$735$}{$637$}{$728$}{$105$}{$686$}{}}
\LANGsk{
\Question{
Súčet všetkých kladných násobkov čísla $7$ menších než $100$ je:}
{$735$}{$637$}{$728$}{$105$}{$686$}{}}
\LANGes{
\Question{
Halla la suma de todos los múltiplos positivos de $7$ menores que $100$.}
{$735$}{$637$}{$728$}{$105$}{$686$}{}}
\End

\Data\ID{1003085204}
\Updated{2020-07-15 03:07:12}
\Level{B}
\SubArea{Arithmetic sequences}
\LANGen{
\Question{
Evaluate the sum: $\sum\limits_{n=3}^{15}(-2n+5)$.}
{$-169$}{$299$}{$-156$}{$-338$}{$-312$}{}}
\LANGpl{
\Question{
Oblicz sumę: $\sum\limits_{n=3}^{15}(-2n+5)$.}
{$-169$}{$299$}{$-156$}{$-338$}{$-312$}{}}
\LANGcs{
\Question{
Vypočítejte součet: $\sum\limits_{n=3}^{15}(-2n+5)$.}
{$-169$}{$299$}{$-156$}{$-338$}{$-312$}{}}
\LANGsk{
\Question{
Vypočítajte súčet: $\sum\limits_{n=3}^{15}(-2n+5)$.}
{$-169$}{$299$}{$-156$}{$-338$}{$-312$}{}}
\LANGes{
\Question{
Halla la suma: $\sum\limits_{n=3}^{15}(-2n+5)$.}
{$-169$}{$299$}{$-156$}{$-338$}{$-312$}{}}
\End

\Data\ID{1003085205}
\Updated{2021-05-03 08:05:40}
\Level{B}
\SubArea{Arithmetic sequences}
\LANGen{
\Question{
Find the sum of the $10$th up to $20$th term of an arithmetic sequence given by: $a_{n+1}=a_n+2$ for all $n\in\mathbb{N}$ and $a_1=-4$.}
{$264$}{$240$}{$480$}{$-320$}{$-352$}{}}
\LANGpl{
\Question{
Wyznacz sumę od  $10$ do $20$ wyrazu ciągu arytmetycznego jeśli : $a_{n+1}=a_n+2$ dla wszystkich  $n\in\mathbb{N}$, a  $a_1=-4$.}
{$264$}{$240$}{$480$}{$-320$}{$-352$}{}}
\LANGcs{
\Question{
Určete součet desátého až dvacátého členu aritmetické posloupnosti, ve které platí: $a_{n+1}=a_n+2$ pro všechna $n\in\mathbb{N}$ a $a_1=-4$.}
{$264$}{$240$}{$480$}{$-320$}{$-352$}{}}
\LANGsk{
\Question{
Určte účet desiateho až dvadsiateho člena aritmetickej postupnosti, v ktorej platí: $a_{n+1}=a_n+2$ pre všetky $n\in\mathbb{N}$ a $a_1=-4$.}
{$264$}{$240$}{$480$}{$-320$}{$-352$}{}}
\LANGes{
\Question{
Halla la suma desde el término $10$ hasta el término $20$ de la progresión aritmética: $a_{n+1}=a_n+2$ para todos $n\in\mathbb{N}$ y $a_1=-4$.}
{$264$}{$240$}{$480$}{$-320$}{$-352$}{}}
\End

\Data\ID{1003085206}
\Updated{2021-05-03 08:05:00}
\Level{B}
\SubArea{Arithmetic sequences}
\LANGen{
\Question{
Find the sum of the first $15$ positive terms of the arithmetic sequence: $\left(3n-7\right)_{n=1}^{\infty}$.}
{$345$}{$255$}{$238$}{$260$}{$322$}{}}
\LANGpl{
\Question{
Wyznacz sumę pierwszych $15$ dodatnich wyrazów ciągu arytmetycznego: $\left\{(3n-7)\right\}_{n=1}^{\infty}$.}
{$345$}{$255$}{$238$}{$260$}{$322$}{}}
\LANGcs{
\Question{
V aritmetické posloupnosti $(3n-7)_{n=1}^{\infty}$ určete součet prvních patnácti kladných členů.}
{$345$}{$255$}{$238$}{$260$}{$322$}{}}
\LANGsk{
\Question{
V aritmetickej postupnosti $(3n-7)_{n=1}^{\infty}$ určte súčet prvých pätnástich kladných členov.}
{$345$}{$255$}{$238$}{$260$}{$322$}{}}
\LANGes{
\Question{
Halla la suma de los $15$ primeros términos positivos de la progresión aritmética: $\left\{(3n-7)\right\}_{n=1}^{\infty}$.}
{$345$}{$255$}{$238$}{$260$}{$322$}{}}
\End

\Data\ID{1003085207}
\Updated{2020-07-15 03:07:12}
\Level{B}
\SubArea{Arithmetic sequences}
\LANGen{
\Question{
Given an arithmetic sequence $\{a_n\}_{n=1}^{\infty}$, where $a_5=\frac{11}2$, $a_{10}=-2$ and $s_k=a_1+\dots+a_k=15$, find $k$.}
{$15$}{$11$}{$10$}{$9$}{$13$}{}}
\LANGpl{
\Question{
Dany jest ciąg arytmetyczny $\{a_n\}_{n=1}^{\infty}$, gdzie $a_5=\frac{11}2$, $a_{10}=-2$ , a $s_k=a_1+\dots+a_k=15$, wyznacz $k$.}
{$15$}{$11$}{$10$}{$9$}{$13$}{}}
\LANGcs{
\Question{
V aritmetické posloupnosti je dáno: $a_5=\frac{11}2$, $a_{10}=-2$ a $s_k=a_1+\dots+a_k=15$. Určete $k$.}
{$15$}{$11$}{$10$}{$9$}{$13$}{}}
\LANGsk{
\Question{
V aritmetickej postupnosti je dané: $a_5=\frac{11}2$, $a_{10}=-2$ a $s_k=a_1+\dots+a_k=15$. Určte $k$.}
{$15$}{$11$}{$10$}{$9$}{$13$}{}}
\LANGes{
\Question{
Dada la progresión aritmética $\{a_n\}_{n=1}^{\infty}$, donde $a_5=\frac{11}2$, $a_{10}=-2$ y $s_k=a_1+\dots+a_k=15$, halla $k$.}
{$15$}{$11$}{$10$}{$9$}{$13$}{}}
\End

\Data\ID{1003085209}
\Updated{2020-07-15 03:07:12}
\Level{B}
\SubArea{Arithmetic sequences}
\LANGen{
\Question{
The sum of the $2$nd and the $4$th term of an arithmetic sequence is $1.2$, the sum of the first $10$ terms is $3.5$. Find the first term.}
{$0.8$}{$0.3$}{$0.35$}{$-0.1$}{$-0.5$}{}}
\LANGpl{
\Question{
Suma  $2$ i  $4$ wyrazu ciągu arytmetycznego wynosi  $1{,}2$, suma pierwszych $10$ wyrazów jest równa  $3{,}5$. Wyznacz pierwszy wyraz tego ciągu.}
{$0{,}8$}{$0{,}3$}{$0{,}35$}{$-0{,}1$}{$-0{,}5$}{}}
\LANGcs{
\Question{
Součet druhého a čtvrtého členu aritmetické posloupnosti je $1{,}2$, součet prvních deseti členů je $3{,}5$. Určete první člen.}
{$0{,}8$}{$0{,}3$}{$0{,}35$}{$-0{,}1$}{$-0{,}5$}{}}
\LANGsk{
\Question{
Súčet druhého a štvrtého člena aritmetickej postupnosti je $1{,}2$, súčet prvých desiatich členov je $3{,}5$. Určte prvý člen.}
{$0{,}8$}{$0{,}3$}{$0{,}35$}{$-0{,}1$}{$-0{,}5$}{}}
\LANGes{
\Question{
La suma del segundo término y del cuarto término de una progresión aritmética es $1.2$, la suma de los $10$ primeros términos es $3.5$. Halla el primer término.}
{$0.8$}{$0.3$}{$0.35$}{$-0.1$}{$-0.5$}{}}
\End

\Data\ID{1003085210}
\Updated{2020-07-15 03:07:12}
\Level{B}
\SubArea{Arithmetic sequences}
\LANGen{
\Question{
The common difference of an arithmetic sequence $\{a_n\}_{n=1}^{\infty}$ is $-\mathrm{e}$. Find the second term, if the sum of the $4$th term up to $9$th term is $6\pi-27\mathrm{e}$.}
{$\pi$}{$\pi+\mathrm{e}$}{$\pi-\mathrm{e}$}{$\pi-2\mathrm{e}$}{$\mathrm{e}-\pi$}{}}
\LANGpl{
\Question{
Wspólna różnica ciągu arytmetycznego  $\{a_n\}_{n=1}^{\infty}$ wynosi $-\mathrm{e}$. Wyznacz drugi wyraz, jeśli suma od $4$-tego  do $9$-tego wyrazu jest równa  $6\pi-27\mathrm{e}$.}
{$\pi$}{$\pi+\mathrm{e}$}{$\pi-\mathrm{e}$}{$\pi-2\mathrm{e}$}{$\mathrm{e}-\pi$}{}}
\LANGcs{
\Question{
Diference aritmetické posloupnosti je $-\mathrm{e}$, součet čtvrtého až devátého členu je $6\pi-27\mathrm{e}$. Určete druhý člen.}
{$\pi$}{$\pi+\mathrm{e}$}{$\pi-\mathrm{e}$}{$\pi-2\mathrm{e}$}{$\mathrm{e}-\pi$}{}}
\LANGsk{
\Question{
Diferencia aritmetickej postupnosti je $-\mathrm{e}$, súčet štvrtého až deviateho člena je $6\pi-27\mathrm{e}$. Určte druhý člen.}
{$\pi$}{$\pi+\mathrm{e}$}{$\pi-\mathrm{e}$}{$\pi-2\mathrm{e}$}{$\mathrm{e}-\pi$}{}}
\LANGes{
\Question{
La diferencia de una progresión aritmética $\{a_n\}_{n=1}^{\infty}$ es $-\mathrm{e}$. Halla el segundo término, si la suma desde el cuarto término hasta el término $9$ es $6\pi-27\mathrm{e}$.}
{$\pi$}{$\pi+\mathrm{e}$}{$\pi-\mathrm{e}$}{$\pi-2\mathrm{e}$}{$\mathrm{e}-\pi$}{}}
\End

\Data\ID{1003097001}
\Updated{2021-05-03 08:05:46}
\Level{B}
\SubArea{Arithmetic sequences}
\LANGen{
\Question{
Find the first term of an arithmetic sequence \( \left\{a_n\right\}_{n=1}^{\infty} \), if:
\[\begin{aligned}
a_2+a_9&=24 \\
a_8-a_6&=4
\end{aligned} \]}
{\( 3 \)}{\( 6 \)}{\( 2 \)}{\( 10 \)}{\( 1 \)}{}}
\LANGpl{
\Question{
Wyznacz pierwszy wyraz ciągu arytmetycznego  \( \left\{a_n\right\}_{n=1}^{\infty} \), jeśli:
\[\begin{aligned}
a_2+a_9&=24 \\
a_8-a_6&=4
\end{aligned} \]}
{\( 3 \)}{\( 6 \)}{\( 2 \)}{\( 10 \)}{\( 1 \)}{}}
\LANGcs{
\Question{
Určete první člen aritmetické posloupnosti \( \left\{a_n\right\}_{n=1}^{\infty} \), platí-li:
\[\begin{aligned}
a_2+a_9&=24 \\
a_8-a_6&=4
\end{aligned} \]}
{\( 3 \)}{\( 6 \)}{\( 2 \)}{\( 10 \)}{\( 1 \)}{}}
\LANGsk{
\Question{
Určte prvý člen aritmetickej postupnosti \( \left\{a_n\right\}_{n=1}^{\infty} \), ak platí:
\[\begin{aligned}
a_2+a_9&=24 \\
a_8-a_6&=4
\end{aligned} \]}
{\( 3 \)}{\( 6 \)}{\( 2 \)}{\( 10 \)}{\( 1 \)}{}}
\LANGes{
\Question{
Halla el primer término de una progresión aritmética \( \left\{a_n\right\}_{n=1}^{\infty} \), si sabemos que:
\[\begin{aligned}
a_2+a_9&=24 \\
a_8-a_6&=4
\end{aligned} \]}
{\( 3 \)}{\( 6 \)}{\( 2 \)}{\( 10 \)}{\( 1 \)}{}}
\End

\Data\ID{1003097002}
\Updated{2021-05-03 08:05:19}
\Level{B}
\SubArea{Arithmetic sequences}
\LANGen{
\Question{
Find the common difference of an arithmetic sequence \( \left\{a_n\right\}_{n=1}^{\infty} \), if:
\[ \begin{aligned}
a_{11}+a_4&=-19 \\ 
a_6+a_7&=-13
\end{aligned} \]}
{\( -3 \)}{\( 3 \)}{\( 7 \)}{\( -1 \)}{\( 10 \)}{}}
\LANGpl{
\Question{
Wyznacz wspólną różnicę ciągu arytmetycznego \( \left\{a_n\right\}_{n=1}^{\infty} \), jeśli:
\[ \begin{aligned}
a_{11}+a_4&=-19 \\ 
a_6+a_7&=-13
\end{aligned} \]}
{\( -3 \)}{\( 3 \)}{\( 7 \)}{\( -1 \)}{\( 10 \)}{}}
\LANGcs{
\Question{
Určete diferenci aritmetické posloupnosti \( \left\{a_n\right\}_{n=1}^{\infty} \), platí-li:
\[ \begin{aligned}
a_{11}+a_4&=-19 \\ 
a_6+a_7&=-13
\end{aligned} \]}
{\( -3 \)}{\( 3 \)}{\( 7 \)}{\( -1 \)}{\( 10 \)}{}}
\LANGsk{
\Question{
Určte diferenciu aritmetickej postupnosti \( \left\{a_n\right\}_{n=1}^{\infty} \), ak platí:
\[ \begin{aligned}
a_{11}+a_4&=-19 \\ 
a_6+a_7&=-13
\end{aligned} \]}
{\( -3 \)}{\( 3 \)}{\( 7 \)}{\( -1 \)}{\( 10 \)}{}}
\LANGes{
\Question{
Halla la diferencia de una progresión aritmética \( \left\{a_n\right\}_{n=1}^{\infty} \), si sabemos que:
\[ \begin{aligned}
a_{11}+a_4&=-19 \\ 
a_6+a_7&=-13
\end{aligned} \]}
{\( -3 \)}{\( 3 \)}{\( 7 \)}{\( -1 \)}{\( 10 \)}{}}
\End

\Data\ID{1003097003}
\Updated{2021-05-03 08:05:33}
\Level{B}
\SubArea{Arithmetic sequences}
\LANGen{
\Question{
Find the sum of the first five terms of an arithmetic sequence \( \left\{a_n\right\}_{n=1}^{\infty} \), if:
\[ \begin{aligned}
a_1\cdot a_4&=100 \\ 
a_2+a_6&=10
\end{aligned} \]}
{\( 50 \)}{\( 0 \)}{\( 20 \)}{\( -5 \)}{\( 90 \)}{}}
\LANGpl{
\Question{
Wyznacz sumę pięciu pierwszych wyrazów ciągu arytmetycznego  \( \left\{a_n\right\}_{n=1}^{\infty} \), jeśli:
\[ \begin{aligned}
a_1\cdot a_4&=100 \\ 
a_2+a_6&=10
\end{aligned} \]}
{\( 50 \)}{\( 0 \)}{\( 20 \)}{\( -5 \)}{\( 90 \)}{}}
\LANGcs{
\Question{
Určete součet prvních pěti členů aritmetické posloupnosti \( \left\{a_n\right\}_{n=1}^{\infty} \), platí-li:
\[ \begin{aligned}
a_1\cdot a_4&=100 \\ 
a_2+a_6&=10
\end{aligned} \]}
{\( 50 \)}{\( 0 \)}{\( 20 \)}{\( -5 \)}{\( 90 \)}{}}
\LANGsk{
\Question{
Určte súčet prvých piatich členov aritmetickej postupnosti \( \left\{a_n\right\}_{n=1}^{\infty} \), ak platí:
\[ \begin{aligned}
a_1\cdot a_4&=100 \\ 
a_2+a_6&=10
\end{aligned} \]}
{\( 50 \)}{\( 0 \)}{\( 20 \)}{\( -5 \)}{\( 90 \)}{}}
\LANGes{
\Question{
Halla la suma se los cinco primeros términos de una progresión aritmética \( \left\{a_n\right\}_{n=1}^{\infty} \), si sabemos que:
\[ \begin{aligned}
a_1\cdot a_4&=100 \\ 
a_2+a_6&=10
\end{aligned} \]}
{\( 50 \)}{\( 0 \)}{\( 20 \)}{\( -5 \)}{\( 90 \)}{}}
\End

\Data\ID{1003097004}
\Updated{2021-05-03 08:05:45}
\Level{B}
\SubArea{Arithmetic sequences}
\LANGen{
\Question{
Find the third term of an arithmetic sequence \( \left\{a_n\right\}_{n=1}^{\infty} \), if:
\[ \begin{aligned}
2a_{10}-3a_3&=5 \\ 
5a_5+4a_1&=83
\end{aligned} \]}
{\( 9 \)}{\( 7 \)}{\( 1 \)}{\( 3 \)}{\( 11 \)}{}}
\LANGpl{
\Question{
Wyznacz trzeci wyraz ciągu arytmetycznego  \( \left\{a_n\right\}_{n=1}^{\infty} \), jeśli:
\[ \begin{aligned}
2a_{10}-3a_3&=5 \\ 
5a_5+4a_1&=83
\end{aligned} \]}
{\( 9 \)}{\( 7 \)}{\( 1 \)}{\( 3 \)}{\( 11 \)}{}}
\LANGcs{
\Question{
Určete třetí člen aritmetické posloupnosti \( \left\{a_n\right\}_{n=1}^{\infty} \), platí-li:
\[ \begin{aligned}
2a_{10}-3a_3&=5 \\ 
5a_5+4a_1&=83
\end{aligned} \]}
{\( 9 \)}{\( 7 \)}{\( 1 \)}{\( 3 \)}{\( 11 \)}{}}
\LANGsk{
\Question{
Určte tretí člen aritmetickej postupnosti \( \left\{a_n\right\}_{n=1}^{\infty} \), ak platí:
\[ \begin{aligned}
2a_{10}-3a_3&=5 \\ 
5a_5+4a_1&=83
\end{aligned} \]}
{\( 9 \)}{\( 7 \)}{\( 1 \)}{\( 3 \)}{\( 11 \)}{}}
\LANGes{
\Question{
Halla el tercer término de una progresión aritmética \( \left\{a_n\right\}_{n=1}^{\infty} \), si sabemos que:
\[ \begin{aligned}
2a_{10}-3a_3&=5 \\ 
5a_5+4a_1&=83
\end{aligned} \]}
{\( 9 \)}{\( 7 \)}{\( 1 \)}{\( 3 \)}{\( 11 \)}{}}
\End

\Data\ID{1003097005}
\Updated{2021-05-03 08:05:09}
\Level{B}
\SubArea{Arithmetic sequences}
\LANGen{
\Question{
Find the square of the first term of an arithmetic sequence \( \left\{a_n\right\}_{n=1}^{\infty} \), if:
\[ \begin{aligned}
a_2^2+a_3^2&=100 \\
a_5+a_7&=0
\end{aligned} \]}
{\( 100 \)}{\( 196 \)}{\( 64 \)}{\( 169 \)}{\( 200 \)}{}}
\LANGpl{
\Question{
Wyznacz kwadrat pierwszego wyrazu ciągu arytmetycznego  \( \left\{a_n\right\}_{n=1}^{\infty} \), jeśli:
\[ \begin{aligned}
a_2^2+a_3^2&=100 \\
a_5+a_7&=0
\end{aligned} \]}
{\( 100 \)}{\( 196 \)}{\( 64 \)}{\( 169 \)}{\( 200 \)}{}}
\LANGcs{
\Question{
Určete druhou mocninu prvního členu aritmetické posloupnosti \( \left\{a_n\right\}_{n=1}^{\infty} \), platí-li:
\[ \begin{aligned}
a_2^2+a_3^2&=100 \\
a_5+a_7&=0
\end{aligned} \]}
{\( 100 \)}{\( 196 \)}{\( 64 \)}{\( 169 \)}{\( 200 \)}{}}
\LANGsk{
\Question{
Určte druhú mocninu prvého člena aritmetickej postupnosti \( \left\{a_n\right\}_{n=1}^{\infty} \), ak platí:
\[ \begin{aligned}
a_2^2+a_3^2&=100 \\
a_5+a_7&=0
\end{aligned} \]}
{\( 100 \)}{\( 196 \)}{\( 64 \)}{\( 169 \)}{\( 200 \)}{}}
\LANGes{
\Question{
Halla el cuadrado del primer término de una progresión aritmética \( \left\{a_n\right\}_{n=1}^{\infty} \), si sabemos que:
\[ \begin{aligned}
a_2^2+a_3^2&=100 \\
a_5+a_7&=0
\end{aligned} \]}
{\( 100 \)}{\( 196 \)}{\( 64 \)}{\( 169 \)}{\( 200 \)}{}}
\End

\Data\ID{1003097006}
\Updated{2021-05-03 08:05:31}
\Level{B}
\SubArea{Arithmetic sequences}
\LANGen{
\Question{
Find the sum of the first ten terms of an arithmetic sequence \( \left\{a_n\right\}_{n=1}^{\infty} \), if:
\[ \begin{aligned}
a_1+a_5+a_{10}&=40 \\
a_{10}-a_5-a_1&=24
\end{aligned} \]}
{\( 140 \)}{\( 180 \)}{\( 160 \)}{\( 320 \)}{\( 200 \)}{}}
\LANGpl{
\Question{
Wyznacz sumę dziesięciu pierwszych wyrazów ciągu arytmetycznego \( \left\{a_n\right\}_{n=1}^{\infty} \), jeśli:
\[ \begin{aligned}
a_1+a_5+a_{10}&=40 \\
a_{10}-a_5-a_1&=24
\end{aligned} \]}
{\( 140 \)}{\( 180 \)}{\( 160 \)}{\( 320 \)}{\( 200 \)}{}}
\LANGcs{
\Question{
Určete součet prvních deseti členů aritmetické posloupnosti \( \left\{a_n\right\}_{n=1}^{\infty} \), platí-li:
\[ \begin{aligned}
a_1+a_5+a_{10}&=40 \\
a_{10}-a_5-a_1&=24
\end{aligned} \]}
{\( 140 \)}{\( 180 \)}{\( 160 \)}{\( 320 \)}{\( 200 \)}{}}
\LANGsk{
\Question{
Určte súčet prvých desiatich členov aritmetickej postupnosti \( \left\{a_n\right\}_{n=1}^{\infty} \), platí-li:
\[ \begin{aligned}
a_1+a_5+a_{10}&=40 \\
a_{10}-a_5-a_1&=24
\end{aligned} \]}
{\( 140 \)}{\( 180 \)}{\( 160 \)}{\( 320 \)}{\( 200 \)}{}}
\LANGes{
\Question{
Halla la suma de los diez primeros términos de una progresión aritmética \( \left\{a_n\right\}_{n=1}^{\infty} \), si sabemos que:
\[ \begin{aligned}
a_1+a_5+a_{10}&=40 \\
a_{10}-a_5-a_1&=24
\end{aligned} \]}
{\( 140 \)}{\( 180 \)}{\( 160 \)}{\( 320 \)}{\( 200 \)}{}}
\End

\Data\ID{1003097007}
\Updated{2020-07-15 03:07:21}
\Level{B}
\SubArea{Arithmetic sequences}
\LANGen{
\Question{
Find the common difference of an arithmetic sequence \( \left\{a_n\right\}_{n=1}^{\infty} \), if the sum of the first seven terms is $42$ and $a_{10}=-4a_5$:
\[ \begin{aligned}
s_7&=42 \\
a_{10}&=-4a_5
\end{aligned} \]}
{\( -3 \)}{\( 3 \)}{\( 15 \)}{\( 2 \)}{\( -2 \)}{}}
\LANGpl{
\Question{
Wyznacz wspólną różnicę ciągu arytmetycznego  \( \left\{a_n\right\}_{n=1}^{\infty} \), jeśli:
\[ \begin{aligned}
s_7&=42 \\
a_{10}&=-4a_5
\end{aligned} \]}
{\( -3 \)}{\( 3 \)}{\( 15 \)}{\( 2 \)}{\( -2 \)}{}}
\LANGcs{
\Question{
Určete diferenci aritmetické posloupnosti \( \left\{a_n\right\}_{n=1}^{\infty} \), víte-li, že součet prvních sedmi členů je $42$ a platí $a_{10}=-4a_5$:
\[ \begin{aligned}
s_7&=42 \\
a_{10}&=-4a_5
\end{aligned} \]}
{\( -3 \)}{\( 3 \)}{\( 15 \)}{\( 2 \)}{\( -2 \)}{}}
\LANGsk{
\Question{
Určte diferenciu aritmetickej postupnosti \( \left\{a_n\right\}_{n=1}^{\infty} \), viete-li, že súčet prvých siedmich členov je $42$ a platí $a_{10}=-4a_5$:
\[ \begin{aligned}
s_7&=42 \\
a_{10}&=-4a_5
\end{aligned} \]}
{\( -3 \)}{\( 3 \)}{\( 15 \)}{\( 2 \)}{\( -2 \)}{}}
\LANGes{
\Question{
Halla la diferencia de una progresión aritmética \( \left\{a_n\right\}_{n=1}^{\infty} \), si la suma de los siete primeros términos es $42$ y $a_{10}=-4a_5$:
\[ \begin{aligned}
s_7&=42 \\
a_{10}&=-4a_5
\end{aligned} \]}
{\( -3 \)}{\( 3 \)}{\( 15 \)}{\( 2 \)}{\( -2 \)}{}}
\End

\Data\ID{1003097008}
\Updated{2020-11-30 04:11:32}
\Level{B}
\SubArea{Arithmetic sequences}
\LANGen{
\Question{
Find the first term of an arithmetic sequence \( \left\{a_n\right\}_{n=1}^{\infty} \), if the sum of the first eight terms is $12$ and the sum of the first twelve terms is $-6$:
\[ \begin{aligned}
s_8&=12 \\
s_{12}&=-6
\end{aligned} \]}
{\( 5 \)}{\( -5 \)}{\( 2 \)}{\( 1 \)}{\( -3 \)}{}}
\LANGpl{
\Question{
Wyznacz pierwszy argument ciągu arytmetycznego \( \left\{a_n\right\}_{n=1}^{\infty} \), jeśli: 
\[ \begin{aligned}
s_8&=12 \\
s_{12}&=-6
\end{aligned} \]}
{\( 5 \)}{\( -5 \)}{\( 2 \)}{\( 1 \)}{\( -3 \)}{}}
\LANGcs{
\Question{
Určete první člen aritmetické posloupnosti \( \left\{a_n\right\}_{n=1}^{\infty} \), víte-li, že součet prvních osmi členů je $12$ a součet prvních dvanácti členů je $-6$: 
\[ \begin{aligned}
s_8&=12 \\
s_{12}&=-6
\end{aligned} \]}
{\( 5 \)}{\( -5 \)}{\( 2 \)}{\( 1 \)}{\( -3 \)}{}}
\LANGsk{
\Question{
Určte prvý člen aritmetickej postupnosti \( \left\{a_n\right\}_{n=1}^{\infty} \), viete-li, že súčet prvých ôsmich čelných je $ 12 $ a súčet prvých dvanástich členov je $ -6 $
\[ \begin{aligned}
s_8&=12 \\
s_{12}&=-6
\end{aligned} \]}
{\( 5 \)}{\( -5 \)}{\( 2 \)}{\( 1 \)}{\( -3 \)}{}}
\LANGes{
\Question{
Halla el primer término de una progresión aritmética \( \left\{a_n\right\}_{n=1}^{\infty} \), si la suma de los ocho primeros términos es $12$ y la suma de los doce primeros términos es $-6$:
\[ \begin{aligned}
s_8&=12 \\
s_{12}&=-6
\end{aligned} \]}
{\( 5 \)}{\( -5 \)}{\( 2 \)}{\( 1 \)}{\( -3 \)}{}}
\End

\Data\ID{9000060601}
\Updated{2020-07-15 03:07:37}
\Level{B}
\SubArea{Arithmetic sequences}
\LANGen{
\Question{
Identify the real number \(x\) which
ensures that the numbers \(a_{1} = 10^{2}\),
\(a_{2} = 10^{3}\) and
\(a_{3} = x\) are
three consecutive terms of an arithmetic sequence.}
{\(x = 1\: 900\)}{\(x = 1\: 000\)}{\(x = 10\: 000\)}{\(x = 1\: 990\)}{\(x = 100\: 000\)}{}}
\LANGpl{
\Question{
Oblicz wartość liczby rzeczywistej  \(x\), dla której 
liczby  \(a_{1} = 10^{2}\),
\(a_{2} = 10^{3}\) i
\(a_{3} = x\) są trzema kolejnymi wyrazami ciągu arytmetycznego.}
{\(x = 1\: 900\)}{\(x = 1\: 000\)}{\(x = 10\: 000\)}{\(x = 1\: 990\)}{\(x = 100\: 000\)}{}}
\LANGcs{
\Question{
Určete reálné číslo \(x\) 
tak, aby čísla \(a_{1} = 10^{2}\),
\(a_{2} = 10^{3}\),
\(a_{3} = x\) 
tvořila tři po sobě jdoucí členy aritmetické posloupnosti.}
{\(x = 1\: 900\)}{\(x = 1\: 000\)}{\(x = 10\: 000\)}{\(x = 1\: 990\)}{\(x = 100\: 000\)}{}}
\LANGsk{
\Question{
Určte reálne číslo \(x\) 
tak, aby čísla \(a_{1} = 10^{2}\),
\(a_{2} = 10^{3}\),
\(a_{3} = x\) 
tvorili tri po sebe idúce členy aritmetickej postupnosti.}
{\(x = 1\: 900\)}{\(x = 1\: 000\)}{\(x = 10\: 000\)}{\(x = 1\: 990\)}{\(x = 100\: 000\)}{}}
\LANGes{
\Question{
Identifica el número real \(x\) para que los números \(a_{1} = 10^{2}\),
\(a_{2} = 10^{3}\) y
\(a_{3} = x\) sean tres términos consecutivos de una progresión aritmética.}
{\(x = 1\: 900\)}{\(x = 1\: 000\)}{\(x = 10\: 000\)}{\(x = 1\: 990\)}{\(x = 100\: 000\)}{}}
\End

\Data\ID{9000060602}
\Updated{2020-07-15 03:07:37}
\Level{B}
\SubArea{Arithmetic sequences}
\LANGen{
\Question{
Identify the real number \(x\) which
ensures that the numbers \(a_{1} = -12\),
\(a_{2} = x\) and
\(a_{3} = 24\) are
three consecutive terms of an arithmetic sequence.}
{\(x = 6\)}{\(x = 0\)}{\(x = 12\)}{\(x = -24\)}{\(x = -2\)}{}}
\LANGpl{
\Question{
Oblicz wartość liczby rzeczywistej  \(x\), dla której liczby  \(a_{1} = -12\),
\(a_{2} = x\) i 
\(a_{3} = 24\) są trzema kolejnymi wyrazami ciągu arytmetycznego.}
{\(x = 6\)}{\(x = 0\)}{\(x = 12\)}{\(x = -24\)}{\(x = -2\)}{}}
\LANGcs{
\Question{
Určete reálné číslo \(x\) 
tak, aby čísla \(a_{1} = -12\),
\(a_{2} = x\),
\(a_{3} = 24\) 
tvořila tři po sobě jdoucí členy aritmetické posloupnosti.}
{\(x = 6\)}{\(x = 0\)}{\(x = 12\)}{\(x = -24\)}{\(x = -2\)}{}}
\LANGsk{
\Question{
Určte reálne číslo \(x\) 
tak, aby čísla \(a_{1} = -12\),
\(a_{2} = x\),
\(a_{3} = 24\) 
tvorili tri po sebe idúce členy aritmetickej postupnosti.}
{\(x = 6\)}{\(x = 0\)}{\(x = 12\)}{\(x = -24\)}{\(x = -2\)}{}}
\LANGes{
\Question{
Identifica el número real \(x\) para que los números \(a_{1} = -12\),
\(a_{2} = x\) y
\(a_{3} = 24\) sean tres términos consecutivos de una progresión aritmética.}
{\(x = 6\)}{\(x = 0\)}{\(x = 12\)}{\(x = -24\)}{\(x = -2\)}{}}
\End

\Data\ID{9000060603}
\Updated{2020-09-17 01:09:58}
\Level{B}
\SubArea{Arithmetic sequences}
\LANGen{
\Question{
Identify the real number \(x\) which
ensures that the numbers \(a_{1} = -x\),
\(a_{2} = -5\) and
\(a_{3} = 0\) are
three consecutive terms of an arithmetic sequence.}
{\(x = 10\)}{\(x = 0\)}{\(x = -5\)}{\(x = 5\)}{\(x = -10\)}{}}
\LANGpl{
\Question{
Oblicz wartość liczby rzeczywistej  \(x\), dla której liczby  \(a_{1} = -x\),
\(a_{2} = -5\) i
\(a_{3} = 0\) są trzema kolejnymi wyrazami ciągu arytmetycznego.}
{\(x = 10\)}{\(x = 0\)}{\(x = -5\)}{\(x = 5\)}{\(x = -10\)}{}}
\LANGcs{
\Question{
Určete reálné číslo \(x\) 
tak, aby čísla \(a_{1} = -x\),
\(a_{2} = -5\),
\(a_{3} = 0\) 
tvořila tři po sobě jdoucí členy aritmetické posloupnosti.}
{\(x = 10\)}{\(x = 0\)}{\(x = -5\)}{\(x = 5\)}{\(x = -10\)}{}}
\LANGsk{
\Question{
Určte reálne číslo \(x\) 
tak, aby čísla \(a_{1} = -x\),
\(a_{2} = -5\),
\(a_{3} = 0\) 
tvorili tri po sebe idúce členy aritmetickej postupnosti.}
{\(x = 10\)}{\(x = 0\)}{\(x = -5\)}{\(x = 5\)}{\(x = -10\)}{}}
\LANGes{
\Question{
Identifica el número real \(x\) para que los números \(a_{1} = -x\),
\(a_{2} = -5\) y
\(a_{3} = 0\) sean tres términos consecutivos de una progresión aritmética.}
{\(x = 10\)}{\(x = 0\)}{\(x = -5\)}{\(x = 5\)}{\(x = -10\)}{}}
\End

\Data\ID{9000060604}
\Updated{2020-07-15 03:07:37}
\Level{B}
\SubArea{Arithmetic sequences}
\LANGen{
\Question{
Identify the real number \(x\) which
ensures that the numbers \(a_{1} = x\),
\(a_{2} = x + 2\) and
\(a_{3} = 2x\) are
three consecutive terms of an arithmetic sequence.}
{\(x = 4\)}{\(x = 0\)}{\(x = 2\)}{\(x = 6\)}{\(x = 8\)}{}}
\LANGpl{
\Question{
Oblicz wartość liczby rzeczywistej  \(x\), dla której liczby \(a_{1} = x\),
\(a_{2} = x + 2\) i
\(a_{3} = 2x\) są trzema kolejnymi wyrazami ciągu arytmetycznego.}
{\(x = 4\)}{\(x = 0\)}{\(x = 2\)}{\(x = 6\)}{\(x = 8\)}{}}
\LANGcs{
\Question{
Určete reálné číslo \(x\) 
tak, aby čísla \(a_{1} = x\),
\(a_{2} = x + 2\),
\(a_{3} = 2x\) 
tvořila tři po sobě jdoucí členy aritmetické posloupnosti.}
{\(x = 4\)}{\(x = 0\)}{\(x = 2\)}{\(x = 6\)}{\(x = 8\)}{}}
\LANGsk{
\Question{
Určte reálne číslo \(x\) 
tak, aby čísla \(a_{1} = x\),
\(a_{2} = x + 2\),
\(a_{3} = 2x\) 
tvorili tri po sebe idúce členy aritmetickej postupnosti.}
{\(x = 4\)}{\(x = 0\)}{\(x = 2\)}{\(x = 6\)}{\(x = 8\)}{}}
\LANGes{
\Question{
Identifica el número real \(x\) para que los números \(a_{1} = x\),
\(a_{2} = x + 2\) y
\(a_{3} = 2x\) sean tres términos consecutivos de una progresión aritmética.}
{\(x = 4\)}{\(x = 0\)}{\(x = 2\)}{\(x = 6\)}{\(x = 8\)}{}}
\End

\Data\ID{9000060605}
\Updated{2020-07-15 03:07:37}
\Level{B}
\SubArea{Arithmetic sequences}
\LANGen{
\Question{
Identify the real number \(x\) which
ensures that the numbers \(a_{1} = x^{2} + 10\),
\(a_{2} = x^{2} + 2x\) and
\(a_{3} = x^{2}\) are
three consecutive terms of an arithmetic sequence.}
{\(x = 2.5\)}{\(x = 0\)}{\(x = 2\)}{\(x = 5\)}{\(x = -5\)}{}}
\LANGpl{
\Question{
Oblicz wartość liczby rzeczywistej  \(x\), dla której liczby  \(a_{1} = x^{2} + 10\),
\(a_{2} = x^{2} + 2x\) i
\(a_{3} = x^{2}\) są trzema kolejnymi wyrazami ciągu arytmetycznego.}
{\(x = 2.5\)}{\(x = 0\)}{\(x = 2\)}{\(x = 5\)}{\(x = -5\)}{}}
\LANGcs{
\Question{
Určete reálné číslo \(x\) 
tak, aby čísla \(a_{1} = x^{2} + 10\),
\(a_{2} = x^{2} + 2x\),
\(a_{3} = x^{2}\) 
tvořila tři po sobě jdoucí členy aritmetické posloupnosti.}
{\(x = 2{,}5\)}{\(x = 0\)}{\(x = 2\)}{\(x = 5\)}{\(x = -5\)}{}}
\LANGsk{
\Question{
Určte reálne číslo \(x\) 
tak, aby čísla \(a_{1} = x^{2} + 10\),
\(a_{2} = x^{2} + 2x\),
\(a_{3} = x^{2}\) 
tvorili tri po sebe idúce členy aritmetickej postupnosti.}
{\(x = 2{,}5\)}{\(x = 0\)}{\(x = 2\)}{\(x = 5\)}{\(x = -5\)}{}}
\LANGes{
\Question{
Identifica el número real \(x\) para que los números \(a_{1} = x^{2} + 10\),
\(a_{2} = x^{2} + 2x\) y
\(a_{3} = x^{2}\) sean tres términos consecutivos de una progresión aritmética.}
{\(x = 2.5\)}{\(x = 0\)}{\(x = 2\)}{\(x = 5\)}{\(x = -5\)}{}}
\End

\Data\ID{9000060606}
\Updated{2020-09-17 01:09:42}
\Level{B}
\SubArea{Arithmetic sequences}
\LANGen{
\Question{
Identify the real number \(x\) which
ensures that the numbers \(a_{1} = 5x + 1\),
\(a_{2} = x\) and
\(a_{3} = 7x + 3\) are
three consecutive terms of an arithmetic sequence.}
{\(x = -0.4\)}{\(x = 0.4\)}{\(x = 2.5\)}{\(x = -2.5\)}{\(x = 5\)}{}}
\LANGpl{
\Question{
Oblicz wartość liczby rzeczywistej \(x\), dla której liczby \(a_{1} = 5x + 1\),
\(a_{2} = x\) i
\(a_{3} = 7x + 3\) są trzema kolejnymi wyrazami ciągu arytmetycznego.}
{\(x = -0.4\)}{\(x = 0.4\)}{\(x = 2.5\)}{\(x = -2.5\)}{\(x = 5\)}{}}
\LANGcs{
\Question{
Určete reálné číslo \(x\) 
tak, aby čísla \(a_{1} = 5x + 1\),
\(a_{2} = x\),
\(a_{3} = 7x + 3\) 
tvořila tři po sobě jdoucí členy aritmetické posloupnosti.}
{\(x = -0{,}4\)}{\(x = 0{,}4\)}{\(x = 2{,}5\)}{\(x = -2{,}5\)}{\(x = 5\)}{}}
\LANGsk{
\Question{
Určte reálne číslo \(x\) 
tak, aby čísla \(a_{1} = 5x + 1\),
\(a_{2} = x\),
\(a_{3} = 7x + 3\) 
tvorili tri po sebe idúce členy aritmetickej postupnosti.}
{\(x = -0{,}4\)}{\(x = 0{,}4\)}{\(x = 2{,}5\)}{\(x = -2{,}5\)}{\(x = 5\)}{}}
\LANGes{
\Question{
Identifica el número real  \(x\) para que los números \(a_{1} = 5x + 1\),
\(a_{2} = x\) y
\(a_{3} = 7x + 3\) sean tres términos consecutivos de una progresión aritmética.}
{\(x = -0.4\)}{\(x = 0.4\)}{\(x = 2.5\)}{\(x = -2.5\)}{\(x = 5\)}{}}
\End

\Data\ID{9000060607}
\Updated{2020-07-15 03:07:37}
\Level{B}
\SubArea{Arithmetic sequences}
\LANGen{
\Question{
Identify the real number \(x\) which
ensures that the numbers \(a_{1} = x^{2} + 2x\),
\(a_{2} = 2x^{2} + 4x\) and
\(a_{3} = x^{2} - 2x - 8\) are
three consecutive terms of an arithmetic sequence.}
{\(x = -2\)}{\(x = 0\)}{\(x = 2\)}{\(x = 4\)}{\(x = -4\)}{}}
\LANGpl{
\Question{
Oblicz wartość liczby rzeczywistej \(x\), dla której liczby  \(a_{1} = x^{2} + 2x\),
\(a_{2} = 2x^{2} + 4x\) i
\(a_{3} = x^{2} - 2x - 8\) są trzema kolejnymi wyrazami ciągu arytmetycznego.}
{\(x = -2\)}{\(x = 0\)}{\(x = 2\)}{\(x = 4\)}{\(x = -4\)}{}}
\LANGcs{
\Question{
Určete reálné číslo \(x\) 
tak, aby čísla \(a_{1} = x^{2} + 2x\),
\(a_{2} = 2x^{2} + 4x\),
\(a_{3} = x^{2} - 2x - 8\) 
tvořila tři po sobě jdoucí členy aritmetické posloupnosti.}
{\(x = -2\)}{\(x = 0\)}{\(x = 2\)}{\(x = 4\)}{\(x = -4\)}{}}
\LANGsk{
\Question{
Určte reálne číslo \(x\) 
tak, aby čísla \(a_{1} = x^{2} + 2x\),
\(a_{2} = 2x^{2} + 4x\),
\(a_{3} = x^{2} - 2x - 8\) 
tvorili tri po sebe idúce členy aritmetickej postupnosti.}
{\(x = -2\)}{\(x = 0\)}{\(x = 2\)}{\(x = 4\)}{\(x = -4\)}{}}
\LANGes{
\Question{
Identifica el número real \(x\) para que los números \(a_{1} = x^{2} + 2x\),
\(a_{2} = 2x^{2} + 4x\) y
\(a_{3} = x^{2} - 2x - 8\) sean tres términos consecutivos de una progresión aritmética.}
{\(x = -2\)}{\(x = 0\)}{\(x = 2\)}{\(x = 4\)}{\(x = -4\)}{}}
\End

\Data\ID{9000060608}
\Updated{2020-07-15 03:07:37}
\Level{B}
\SubArea{Arithmetic sequences}
\LANGen{
\Question{
Identify the real number \(x\) which
ensures that the numbers \(a_{1} =\log x\),
\(a_{2} =\log(2x)\) and
\(a_{3} = 1\) are
three consecutive terms of an arithmetic sequence.}
{\(x = 2.5\)}{\(x = 1\)}{\(x =\log 2\)}{\(x = 2\)}{\(x = 1\)}{}}
\LANGpl{
\Question{
Oblicz wartość liczby rzeczywistej  \(x\), dla której  \(a_{1} =\log x\),
\(a_{2} =\log(2x)\) i
\(a_{3} = 1\) są trzema kolejnymi wyrazami ciągu arytmetycznego.}
{\(x = 2.5\)}{\(x = 1\)}{\(x =\log 2\)}{\(x = 2\)}{\(x = 1\)}{}}
\LANGcs{
\Question{
Určete reálné číslo \(x\) 
tak, aby čísla \(a_{1} =\log x\),
\(a_{2} =\log(2x)\),
\(a_{3} = 1\) 
tvořila tři po sobě jdoucí členy aritmetické posloupnosti.}
{\(x = 2{,}5\)}{\(x = 1\)}{\(x =\log 2\)}{\(x = 2\)}{\(x = 1\)}{}}
\LANGsk{
\Question{
Určte reálne číslo \(x\) 
tak, aby čísla \(a_{1} =\log x\),
\(a_{2} =\log(2x)\),
\(a_{3} = 1\) 
tvorili tri po sebe idúce členy aritmetickej postupnosti.}
{\(x = 2{,}5\)}{\(x = 1\)}{\(x =\log 2\)}{\(x = 2\)}{\(x = 1\)}{}}
\LANGes{
\Question{
Identifica el número real \(x\) para que los números \(a_{1} =\log x\),
\(a_{2} =\log(2x)\) y
\(a_{3} = 1\) sean tres términos consecutivos de una progresión aritmética.}
{\(x = 2.5\)}{\(x = 1\)}{\(x =\log 2\)}{\(x = 2\)}{\(x = 1\)}{}}
\End

\Data\ID{9000060609}
\Updated{2020-07-15 03:07:37}
\Level{B}
\SubArea{Arithmetic sequences}
\LANGen{
\Question{
Identify the real number \(x\) which
ensures that the numbers \(a_{1} =\log (x + 2)\),
\(a_{2} =\log (3x + 6)\) and
\(a_{3} =\log 18\) are
three consecutive terms of an arithmetic sequence.}
{\(x = 0\)}{\(x =\log 3\)}{\(x = 3\)}{\(x = 8\)}{\(x = 18\)}{}}
\LANGpl{
\Question{
Oblicz wartość liczby rzeczywistej  \(x\), dla której liczby \(a_{1} =\log (x + 2)\),
\(a_{2} =\log (3x + 6)\) i
\(a_{3} =\log 18\) są trzema kolejnymi wyrazami ciągu arytmetycznego.}
{\(x = 0\)}{\(x =\log 3\)}{\(x = 3\)}{\(x = 8\)}{\(x = 18\)}{}}
\LANGcs{
\Question{
Určete reálné číslo \(x\) 
tak, aby čísla \(a_{1} =\log (x + 2)\),
\(a_{2} =\log (3x + 6)\),
\(a_{3} =\log 18\) 
tvořila tři po sobě jdoucí členy aritmetické posloupnosti.}
{\(x = 0\)}{\(x =\log 3\)}{\(x = 3\)}{\(x = 8\)}{\(x = 18\)}{}}
\LANGsk{
\Question{
Určte reálne číslo \(x\) 
tak, aby čísla \(a_{1} =\log (x + 2)\),
\(a_{2} =\log (3x + 6)\),
\(a_{3} =\log 18\) 
tvorili tri po sebe idúce členy aritmetickej postupnosti.}
{\(x = 0\)}{\(x =\log 3\)}{\(x = 3\)}{\(x = 8\)}{\(x = 18\)}{}}
\LANGes{
\Question{
Identifica el número real \(x\) para que los números \(a_{1} =\log (x + 2)\),
\(a_{2} =\log (3x + 6)\) y
\(a_{3} =\log 18\) sean tres términos consecutivos de una progresión aritmética.}
{\(x = 0\)}{\(x =\log 3\)}{\(x = 3\)}{\(x = 8\)}{\(x = 18\)}{}}
\End

\Data\ID{9000060610}
\Updated{2020-09-17 01:09:14}
\Level{B}
\SubArea{Arithmetic sequences}
\LANGen{
\Question{
Identify the real number \(x\) which
ensures that the numbers \(a_{1} =\log x\),
\(a_{2} = 2\) and
\(a_{3} =\log x^{3}\) are
three consecutive terms of an arithmetic sequence.}
{\(x = 10\)}{\(x = 0\)}{\(x = 0.1\)}{\(x = 1\)}{\(x = -1\)}{}}
\LANGpl{
\Question{
Oblicz wartość liczby rzeczywistej  \(x\), dla której liczby \(a_{1} =\log x\),
\(a_{2} = 2\) i
\(a_{3} =\log x^{3}\) są trzema kolejnymi wyrazami ciągu arytmetycznego.}
{\(x = 10\)}{\(x = 0\)}{\(x = 0.1\)}{\(x = 1\)}{\(x = -1\)}{}}
\LANGcs{
\Question{
Určete reálné číslo \(x\) 
tak, aby čísla \(a_{1} =\log x\),
\(a_{2} = 2\),
\(a_{3} =\log x^{3}\) 
tvořila tři po sobě jdoucí členy aritmetické posloupnosti.}
{\(x = 10\)}{\(x = 0\)}{\(x = 0{,}1\)}{\(x = 1\)}{\(x = -1\)}{}}
\LANGsk{
\Question{
Určte reálne číslo \(x\) 
tak, aby čísla \(a_{1} =\log x\),
\(a_{2} = 2\),
\(a_{3} =\log x^{3}\) 
tvorili tri po sebe idúce členy aritmetickej postupnosti.}
{\(x = 10\)}{\(x = 0\)}{\(x = 0{,}1\)}{\(x = 1\)}{\(x = -1\)}{}}
\LANGes{
\Question{
Identifica el número real \(x\) para que los números \(a_{1} =\log x\),
\(a_{2} = 2\) y
\(a_{3} =\log x^{3}\) sean tres términos consecutivos de una progresión aritmética.}
{\(x = 10\)}{\(x = 0\)}{\(x = 0.1\)}{\(x = 1\)}{\(x = -1\)}{}}
\End

\Data\ID{9000064806}
\Updated{2021-05-03 08:05:56}
\Level{B}
\SubArea{Arithmetic sequences}
\LANGen{
\Question{
The arithmetic sequence is defined by the first term
\(a_{1} = 17\) and the
fifth term \(a_{5} = 11\).
Find the term which is seven times smaller than the third term of the sequence.}
{\(a_{11}\)}{\(a_{2}\)}{\(a_{8}\)}{\(a_{17}\)}{\(a_{21}\)}{}}
\LANGpl{
\Question{
Dany jest ciąg arytmetyczny, gdzie pierwszy wyraz jest równy 
\(a_{1} = 17\), a piąty 
 \(a_{5} = 11\).
Oblicz wyraz, który jest siedem razy mniejszy od trzeciego wyrazu tego ciągu.}
{\(a_{11}\)}{\(a_{2}\)}{\(a_{8}\)}{\(a_{17}\)}{\(a_{21}\)}{}}
\LANGcs{
\Question{
V aritmetické posloupnosti platí, že
\(a_{1} = 17\),
\(a_{5} = 11\).
Vypočtěte, který člen posloupnosti je sedminou třetího členu.}
{\(a_{11}\)}{\(a_{2}\)}{\(a_{8}\)}{\(a_{17}\)}{\(a_{21}\)}{}}
\LANGsk{
\Question{
V aritmetickej postupnosti platí, že
\(a_{1} = 17\),
\(a_{5} = 11\).
Vypočítajte, ktorý člen postupnosti je sedminou tretieho člena.}
{\(a_{11}\)}{\(a_{2}\)}{\(a_{8}\)}{\(a_{17}\)}{\(a_{21}\)}{}}
\LANGes{
\Question{
Una progresión aritmética viene dada por el primer término
\(a_{1} = 17\) y el quinto término \(a_{5} = 11\).
Halla el término que sea siete veces más pequeño que el tercer término de la progresión.}
{\(a_{11}\)}{\(a_{2}\)}{\(a_{8}\)}{\(a_{17}\)}{\(a_{21}\)}{}}
\End

\Data\ID{9000065307}
\Updated{2020-07-15 03:07:37}
\Level{B}
\SubArea{Arithmetic sequences}
\LANGen{
\Question{
In the arithmetic sequence given by the first term
\(a_{1} = 4\) and the common
difference \(d = 2\) 
find the sum of the first twelve terms of the sequence.}
{\(s_{12} = 180\)}{\(s_{12} = 72\)}{\(s_{12} = 120\)}{\(s_{12} = 168\)}{}{}}
\LANGpl{
\Question{
Dany jest ciąg arytmetyczny, gdzie pierwszy wyraz jest równy 
\(a_{1} = 4\), a różnica  \(d = 2\). Oblicz sumę dwunastu początkowych wyrazów tego ciągu.}
{\(s_{12} = 180\)}{\(s_{12} = 72\)}{\(s_{12} = 120\)}{\(s_{12} = 168\)}{}{}}
\LANGcs{
\Question{
Určete součet prvních dvanácti členů aritmetické posloupnosti, je-li dáno
\(a_{1} = 4\),
\(d = 2\).}
{\(s_{12} = 180\)}{\(s_{12} = 72\)}{\(s_{12} = 120\)}{\(s_{12} = 168\)}{}{}}
\LANGsk{
\Question{
Určte súčet prvých dvanástich členov aritmetickej postupnosti, ak je dané
\(a_{1} = 4\),
\(d = 2\).}
{\(s_{12} = 180\)}{\(s_{12} = 72\)}{\(s_{12} = 120\)}{\(s_{12} = 168\)}{}{}}
\LANGes{
\Question{
Una progresión aritmética viene dada por el primer término
\(a_{1} = 4\) y la diferencia \(d = 2\). Halla la suma de los doce primeros términos.}
{\(s_{12} = 180\)}{\(s_{12} = 72\)}{\(s_{12} = 120\)}{\(s_{12} = 168\)}{}{}}
\End

\Data\ID{9000065308}
\Updated{2020-07-15 03:07:37}
\Level{B}
\SubArea{Arithmetic sequences}
\LANGen{
\Question{
For the arithmetic sequence with the first term $a_1=3$ holds $a_n=27$ and the sum of the first $n$ terms is $195$. Find $n$.}
{\(n = 13\)}{\(n = 14\)}{\(n = 15\)}{\(n = 16\)}{}{}}
\LANGpl{
\Question{
Dany jest ciąg arytmetyczny, gdzie 
\(a_{1} = 3\),
\(a_{n} = 27\) i 
\(s_{n} = 195\) w pewnym konkretnym składniku
\(n\).
Oblicz  \(n\).}
{\(n = 13\)}{\(n = 14\)}{\(n = 15\)}{\(n = 16\)}{}{}}
\LANGcs{
\Question{
V aritmetické posloupnosti je dáno \(a_{1} = 3\),
\(a_{n} = 27\),
\(s_{n} = 195\). Určete
číslo \(n\).}
{\(n = 13\)}{\(n = 14\)}{\(n = 15\)}{\(n = 16\)}{}{}}
\LANGsk{
\Question{
V aritmetickej postupnosti je dané \(a_{1} = 3\),
\(a_{n} = 27\),
\(s_{n} = 195\). Určte
číslo \(n\).}
{\(n = 13\)}{\(n = 14\)}{\(n = 15\)}{\(n = 16\)}{}{}}
\LANGes{
\Question{
Una progresión aritmética viene dada por $a_1=3$, $a_n=27$ y la suma de los $n$ primeros términos es $195$. Halla $n$.}
{\(n = 13\)}{\(n = 14\)}{\(n = 15\)}{\(n = 16\)}{}{}}
\End

\Data\ID{9000065310}
\Updated{2021-05-03 08:05:33}
\Level{B}
\SubArea{Arithmetic sequences}
\LANGen{
\Question{
For the arithmetic sequence with \(a_{4} = 11\) 
and \(a_{9} = -24\) 
find the sum of the first fourteen terms.}
{\(- 189\)}{\(189\)}{\(198\)}{\(- 198\)}{}{}}
\LANGpl{
\Question{
Dany jest ciąg arytmetyczny, gdzie \(a_{4} = 11\) 
i \(a_{9} = -24\). Oblicz sumę czternastu początkowych wyrazów tego ciągu.}
{\(- 189\)}{\(189\)}{\(198\)}{\(- 198\)}{}{}}
\LANGcs{
\Question{
Určete součet prvních čtrnácti členů aritmetické posloupnosti, je-li dáno
\(a_{4} = 11\),
\(a_{9} = -24\).}
{\(- 189\)}{\(189\)}{\(198\)}{\(- 198\)}{}{}}
\LANGsk{
\Question{
Určte súčet prvých štrnástich členov aritmetickej postupnosti, ak je dané
\(a_{4} = 11\),
\(a_{9} = -24\).}
{\(- 189\)}{\(189\)}{\(198\)}{\(- 198\)}{}{}}
\LANGes{
\Question{
Suponiendo que \(a_{4} = 11\) 
y \(a_{9} = -24\), halla la suma de los catorce primeros términos de la progresión aritmética.}
{\(- 189\)}{\(189\)}{\(198\)}{\(- 198\)}{}{}}
\End

\Data\ID{9000072701}
\Updated{2020-07-15 03:07:37}
\Level{B}
\SubArea{Arithmetic sequences}
\LANGen{
\Question{
The following numbers form an arithmetic sequence. Find
\(x\).
\[ 
 1\, ,\ x\, ,\ 3
\]}
{\(x = 2\)}{\(x = -2\)}{\(x = 2.5\)}{\(x = 1.5\)}{}{}}
\LANGpl{
\Question{
Podane liczby tworzą ciąg arytmetyczny. Oblicz wartość 
\(x\).
\[ 
 1\, ,\ x\, ,\ 3
\]}
{\(x = 2\)}{\(x = -2\)}{\(x = 2.5\)}{\(x = 1.5\)}{}{}}
\LANGcs{
\Question{
Je dán výčet několika po sobě jdoucích členů
aritmetické posloupnosti. Určete
\(x\).
\[ 
 1\, ,\ x\, ,\ 3
\]}
{\(x = 2\)}{\(x = -2\)}{\(x = 2{,}5\)}{\(x = 1{,}5\)}{}{}}
\LANGsk{
\Question{
Sú dané tri po sebe idúce členy aritmetickej postupnosti.  Určte
 \(x\).
\[ 
 1\, ,\ x\, ,\ 3
\]}
{\(x = 2\)}{\(x = -2\)}{\(x = 2{,}5\)}{\(x = 1{,}5\)}{}{}}
\LANGes{
\Question{
Los siguientes números forman una progresión aritmética. Halla
\(x\).
\[ 
 1\, ,\ x\, ,\ 3
\]}
{\(x = 2\)}{\(x = -2\)}{\(x = 2.5\)}{\(x = 1.5\)}{}{}}
\End

\Data\ID{9000072702}
\Updated{2020-07-15 03:07:37}
\Level{B}
\SubArea{Arithmetic sequences}
\LANGen{
\Question{
The following numbers form an arithmetic sequence. Find
\(x\).
\[ 
 10\, ,\ 20\, ,\ x
\]}
{\(x = 30\)}{\(x = 40\)}{\(x = -20\)}{\(x = -10\)}{}{}}
\LANGpl{
\Question{
Podane liczby tworzą ciąg arytmetyczny. Oblicz wartość 
\(x\).
\[ 
 10\, ,\ 20\, ,\ x
\]}
{\(x = 30\)}{\(x = 40\)}{\(x = -20\)}{\(x = -10\)}{}{}}
\LANGcs{
\Question{
Je dán výčet několika po sobě jdoucích členů aritmetické posloupnosti.
Určete \(x\).
\[ 
 10\, ,\ 20\, ,\ x
\]}
{\(x = 30\)}{\(x = 40\)}{\(x = -20\)}{\(x = -10\)}{}{}}
\LANGsk{
\Question{
Sú dané tri po sebe idúce členy aritmetickej postupnosti. Určte
\(x\).
\[ 
 10\, ,\ 20\, ,\ x
\]}
{\(x = 30\)}{\(x = 40\)}{\(x = -20\)}{\(x = -10\)}{}{}}
\LANGes{
\Question{
Los siguientes números forman una progresión aritmética. Halla
\(x\).
\[ 
 10\, ,\ 20\, ,\ x
\]}
{\(x = 30\)}{\(x = 40\)}{\(x = -20\)}{\(x = -10\)}{}{}}
\End

\Data\ID{9000072703}
\Updated{2020-09-17 01:09:58}
\Level{B}
\SubArea{Arithmetic sequences}
\LANGen{
\Question{
The following numbers form an arithmetic sequence. Find
\(x\).
\[ 
 x\, ,\ 10\, ,\ 5
\]}
{\(x = 15\)}{\(x = 20\)}{\(x = 50\)}{\(x = 5\)}{}{}}
\LANGpl{
\Question{
Podane liczby tworzą ciąg arytmetyczny. Oblicz wartość 
\(x\).
\[ 
 x\, ,\ 10\, ,\ 5
\]}
{\(x = 15\)}{\(x = 20\)}{\(x = 50\)}{\(x = 5\)}{}{}}
\LANGcs{
\Question{
Je dán výčet několika po sobě jdoucích členů aritmetické posloupnosti.
Určete \(x\).
\[ 
 x\, ,\ 10\, ,\ 5
\]}
{\(x = 15\)}{\(x = 20\)}{\(x = 50\)}{\(x = 5\)}{}{}}
\LANGsk{
\Question{
Sú dané tri po sebe idúce členy aritmetickej postupnosti. Určte
\(x\).
\[ 
 x\, ,\ 10\, ,\ 5
\]}
{\(x = 15\)}{\(x = 20\)}{\(x = 50\)}{\(x = 5\)}{}{}}
\LANGes{
\Question{
Los siguientes números forman una progresión aritmética. Halla
\(x\).
\[ 
 x\, ,\ 10\, ,\ 5
\]}
{\(x = 15\)}{\(x = 20\)}{\(x = 50\)}{\(x = 5\)}{}{}}
\End

\Data\ID{9000072704}
\Updated{2020-07-15 03:07:37}
\Level{B}
\SubArea{Arithmetic sequences}
\LANGen{
\Question{
We are given five consecutive terms of an arithmetic sequence. Find
\(x\).
\[ 
 4\, ,\ a\, ,\ 8\, ,\ b\, ,\ x
\]}
{\(x = 12\)}{\(x = 10\)}{\(x = 14\)}{\(x = 16\)}{}{}}
\LANGpl{
\Question{
Podane liczby tworzą ciąg arytmetyczny. Oblicz wartość 
\(x\).
\[ 
 4\, ,\ a\, ,\ 8\, ,\ b\, ,\ x
\]}
{\(x = 12\)}{\(x = 10\)}{\(x = 14\)}{\(x = 16\)}{}{}}
\LANGcs{
\Question{
Je dán výčet několika po sobě jdoucích členů aritmetické posloupnosti.
Určete \(x\).
\[ 
 4\, ,\ a\, ,\ 8\, ,\ b\, ,\ x
\]}
{\(x = 12\)}{\(x = 10\)}{\(x = 14\)}{\(x = 16\)}{}{}}
\LANGsk{
\Question{
Je daných päť po sebe idúcich členov aritmetickej postupnosti. Určte \(x\).
\[ 
 4\, ,\ a\, ,\ 8\, ,\ b\, ,\ x
\]}
{\(x = 12\)}{\(x = 10\)}{\(x = 14\)}{\(x = 16\)}{}{}}
\LANGes{
\Question{
Los siguientes términos consecutivos forman una progresión aritmética. Halla
\(x\).
\[ 
 4\, ,\ a\, ,\ 8\, ,\ b\, ,\ x
\]}
{\(x = 12\)}{\(x = 10\)}{\(x = 14\)}{\(x = 16\)}{}{}}
\End

\Data\ID{9000072705}
\Updated{2020-09-21 02:09:01}
\Level{B}
\SubArea{Arithmetic sequences}
\LANGen{
\Question{
We are given four consecutive terms of an arithmetic sequence. Find
\(x\).
\[ 
 3\, ,\ a\, ,\ 0\, ,\ x
\]}
{\(x = -1.5\)}{\(x = -3\)}{\(x = 6\)}{\(x = -6\)}{}{}}
\LANGpl{
\Question{
Podane liczby tworzą ciąg arytmetyczny. Oblicz wartość 
\(x\).
\[ 
 3\, ,\ a\, ,\ 0\, ,\ x
\]}
{\(x = -1.5\)}{\(x = -3\)}{\(x = 6\)}{\(x = -6\)}{}{}}
\LANGcs{
\Question{
Je dán výčet několika po sobě jdoucích členů aritmetické posloupnosti.
Určete \(x\).
\[ 
 3\, ,\ a\, ,\ 0\, ,\ x
\]}
{\(x = -1{,}5\)}{\(x = -3\)}{\(x = 6\)}{\(x = -6\)}{}{}}
\LANGsk{
\Question{
Je daných päť po sebe idúcich členov aritmetickej postupnosti. Určte
 \(x\).
\[ 
 3\, ,\ a\, ,\ 0\, ,\ x
\]}
{\(x = -1{,}5\)}{\(x = -3\)}{\(x = 6\)}{\(x = -6\)}{}{}}
\LANGes{
\Question{
Los siguientes términos consecutivos forman una progresión aritmética. Halla
\(x\).
\[ 
 3\, ,\ a\, ,\ 0\, ,\ x
\]}
{\(x = -1.5\)}{\(x = -3\)}{\(x = 6\)}{\(x = -6\)}{}{}}
\End

\Data\ID{9000072706}
\Updated{2020-07-15 03:07:37}
\Level{B}
\SubArea{Arithmetic sequences}
\LANGen{
\Question{
We are given five consecutive terms of an arithmetic sequence. Find
\(x\).
\[ 
 5\, ,\ a\, ,\ b\, ,\ x\, ,\ 6
\]}
{\(x = 5.75\)}{\(x = 5.5\)}{\(x = 5.8\)}{\(x = 5\frac{2}
{3}\)}{}{}}
\LANGpl{
\Question{
Podane liczby tworzą ciąg arytmetyczny. Oblicz wartość 
\(x\).
\[ 
 5\, ,\ a\, ,\ b\, ,\ x\, ,\ 6
\]}
{\(x = 5.75\)}{\(x = 5.5\)}{\(x = 5.8\)}{\(x = 5\frac{2}
{3}\)}{}{}}
\LANGcs{
\Question{
Je dán výčet několika po sobě jdoucích členů aritmetické posloupnosti.
Určete \(x\).
\[ 
 5\, ,\ a\, ,\ b\, ,\ x\, ,\ 6
\]}
{\(x = 5{,}75\)}{\(x = 5{,}5\)}{\(x = 5{,}8\)}{\(x = 5\frac{2}
{3}\)}{}{}}
\LANGsk{
\Question{
Je daných päť po sebe idúcich členov aritmetickej postupnosti. Určte 
 \(x\).
\[ 
 5\, ,\ a\, ,\ b\, ,\ x\, ,\ 6
\]}
{\(x = 5{,}75\)}{\(x = 5{,}5\)}{\(x = 5{,}8\)}{\(x = 5\frac{2}
{3}\)}{}{}}
\LANGes{
\Question{
Los siguientes términos consecutivos forman una progresión aritmética. Halla
\(x\).
\[ 
 5\, ,\ a\, ,\ b\, ,\ x\, ,\ 6
\]}
{\(x = 5.75\)}{\(x = 5.5\)}{\(x = 5.8\)}{\(x = 5\frac{2}
{3}\)}{}{}}
\End

\Data\ID{9000072707}
\Updated{2020-07-15 03:07:37}
\Level{B}
\SubArea{Arithmetic sequences}
\LANGen{
\Question{
We are given seven consecutive terms of an arithmetic sequence. Find
\(x\).
\[ 
 100\, ,\ a\, ,\ b\, ,\ x\, ,\ c\, ,\ d\, ,\ 0
\]}
{\(x = 50\)}{\(x = 60\)}{\(x = 40\)}{\(x = 51\)}{}{}}
\LANGpl{
\Question{
Podane liczby tworzą ciąg arytmetyczny. Oblicz wartość 
\(x\).
\[ 
 100\, ,\ a\, ,\ b\, ,\ x\, ,\ c\, ,\ d\, ,\ 0
\]}
{\(x = 50\)}{\(x = 60\)}{\(x = 40\)}{\(x = 51\)}{}{}}
\LANGcs{
\Question{
Je dán výčet několika po sobě jdoucích členů aritmetické posloupnosti.
Určete \(x\).
\[ 
 100\, ,\ a\, ,\ b\, ,\ x\, ,\ c\, ,\ d\, ,\ 0
\]}
{\(x = 50\)}{\(x = 60\)}{\(x = 40\)}{\(x = 51\)}{}{}}
\LANGsk{
\Question{
Je daných sedem po sebe idúcich členov aritmetickej postupnosti. Určte \(x\).
\[ 
 100\, ,\ a\, ,\ b\, ,\ x\, ,\ c\, ,\ d\, ,\ 0
\]}
{\(x = 50\)}{\(x = 60\)}{\(x = 40\)}{\(x = 51\)}{}{}}
\LANGes{
\Question{
Los siguientes términos consecutivos forman una progresión aritmética. Halla
\(x\).
\[ 
 100\, ,\ a\, ,\ b\, ,\ x\, ,\ c\, ,\ d\, ,\ 0
\]}
{\(x = 50\)}{\(x = 60\)}{\(x = 40\)}{\(x = 51\)}{}{}}
\End

\Data\ID{9000072708}
\Updated{2020-09-24 11:09:01}
\Level{B}
\SubArea{Arithmetic sequences}
\LANGen{
\Question{
We are given six consecutive terms of an arithmetic sequence. Find
\(x\).
\[ 
 \frac52,\ a,\ x,\ b,\ c,\ 5
\]}
{\(x = 3.5\)}{\(x = 3\)}{\(x = 4\)}{\(x = 3.75\)}{}{}}
\LANGpl{
\Question{
Podane liczby tworzą ciąg arytmetyczny. Oblicz wartość 
\(x\).
\[ 
 \frac52,\ a,\ x,\ b,\ c,\ 5
\]}
{\(x = 3.5\)}{\(x = 3\)}{\(x = 4\)}{\(x = 3.75\)}{}{}}
\LANGcs{
\Question{
Je dán výčet několika po sobě jdoucích členů aritmetické posloupnosti.
Určete \(x\).
\[ 
 \frac52,\ a,\ x,\ b,\ c,\ 5
\]}
{\(x = 3{,}5\)}{\(x = 3\)}{\(x = 4\)}{\(x = 3{,}75\)}{}{}}
\LANGsk{
\Question{
Je daných šesť po sebe idúcich členov aritmetickej postupnosti. Určte
 \(x\).
\[ 
 \frac52,\ a,\ x,\ b,\ c,\ 5
\]}
{\(x = 3{,}5\)}{\(x = 3\)}{\(x = 4\)}{\(x = 3{,}75\)}{}{}}
\LANGes{
\Question{
Los siguientes términos consecutivos forman una progresión aritmética. Halla
\(x\).
\[ 
 \frac52,\ a,\ x,\ b,\ c,\ 5
\]}
{\(x = 3.5\)}{\(x = 3\)}{\(x = 4\)}{\(x = 3.75\)}{}{}}
\End

\Data\ID{9000072709}
\Updated{2020-07-15 03:07:37}
\Level{B}
\SubArea{Arithmetic sequences}
\LANGen{
\Question{
We are given seven consecutive terms of an arithmetic sequence. Find
\(x\).
\[ 
 x,\ 1,\ a,\ b,\ c,\ d,\ \frac12
\]}
{\(x = 1.1\)}{\(x = 1.5\)}{\(x = -0.5\)}{\(x = 2\)}{}{}}
\LANGpl{
\Question{
Podane liczby tworzą ciąg arytmetyczny. Oblicz wartość 
\(x\).

\[ 
 x,\ 1,\ a,\ b,\ c,\ d,\ \frac12
\]}
{\(x = 1.1\)}{\(x = 1.5\)}{\(x = -0.5\)}{\(x = 2\)}{}{}}
\LANGcs{
\Question{
Je dán výčet několika po sobě jdoucích členů aritmetické posloupnosti.
Určete \(x\).

\[ 
 x,\ 1,\ a,\ b,\ c,\ d,\ \frac12
\]}
{\(x = 1{,}1\)}{\(x = 1{,}5\)}{\(x = -0{,}5\)}{\(x = 2\)}{}{}}
\LANGsk{
\Question{
Je daných sedem po sebe idúcich členov aritmetickej postupnosti. Určte
 \(x\).

\[ 
 x,\ 1,\ a,\ b,\ c,\ d,\ \frac12
\]}
{\(x = 1{,}1\)}{\(x = 1{,}5\)}{\(x = -0{,}5\)}{\(x = 2\)}{}{}}
\LANGes{
\Question{
Los siguientes términos consecutivos forman una progresión aritmética. Halla
\(x\).
\[ 
 x,\ 1,\ a,\ b,\ c,\ d,\ \frac12
\]}
{\(x = 1.1\)}{\(x = 1.5\)}{\(x = -0.5\)}{\(x = 2\)}{}{}}
\End

\Data\ID{9000072710}
\Updated{2020-07-15 03:07:37}
\Level{B}
\SubArea{Arithmetic sequences}
\LANGen{
\Question{
We are given seven consecutive terms of an arithmetic sequence. Find
\(x\).
\[ 
 \frac{4} 
{5}\, ,\ a\, ,\ b\, ,\ 0\, ,\ c\, ,\ d\, ,\ x
\]}
{\(x = -\frac{4}
{5}\)}{\(x = \frac{5} 
{4}\)}{\(x = -\frac{5}
{4}\)}{\(x = -\frac{8}
{5}\)}{}{}}
\LANGpl{
\Question{
Podane wyrazy tworzą ciąg arytmetyczny. Oblicz wartość 
\(x\).
\[ 
 \frac{4} 
{5}\, ,\ a\, ,\ b\, ,\ 0\, ,\ c\, ,\ d\, ,\ x
\]}
{\(x = -\frac{4}
{5}\)}{\(x = \frac{5} 
{4}\)}{\(x = -\frac{5}
{4}\)}{\(x = -\frac{8}
{5}\)}{}{}}
\LANGcs{
\Question{
Je dán výčet několika po sobě jdoucích členů aritmetické posloupnosti.
Určete \(x\).
\[ 
 \frac{4} 
{5}\, ,\ a\, ,\ b\, ,\ 0\, ,\ c\, ,\ d\, ,\ x
\]}
{\(x = -\frac{4}
{5}\)}{\(x = \frac{5} 
{4}\)}{\(x = -\frac{5}
{4}\)}{\(x = -\frac{8}
{5}\)}{}{}}
\LANGsk{
\Question{
Je daných sedem po sebe idúcich členov aritmetickej postupnosti. Určte
 \(x\).
\[ 
 \frac{4} 
{5}\, ,\ a\, ,\ b\, ,\ 0\, ,\ c\, ,\ d\, ,\ x
\]}
{\(x = -\frac{4}
{5}\)}{\(x = \frac{5} 
{4}\)}{\(x = -\frac{5}
{4}\)}{\(x = -\frac{8}
{5}\)}{}{}}
\LANGes{
\Question{
Los siguientes términos consecutivos forman una progresión aritmética. Halla
\(x\).
\[ 
 \frac{4} 
{5}\, ,\ a\, ,\ b\, ,\ 0\, ,\ c\, ,\ d\, ,\ x
\]}
{\(x = -\frac{4}
{5}\)}{\(x = \frac{5} 
{4}\)}{\(x = -\frac{5}
{4}\)}{\(x = -\frac{8}
{5}\)}{}{}}
\End

\Data\ID{1003047801}
\Updated{2021-05-03 07:05:38}
\Level{C}
\SubArea{Arithmetic sequences}
\LANGen{
\Question{
A certain type of bamboo grows $1.3\,\mathrm{m}$ per day during  the vegetation period. Concerning the fact that it reached a height of $30\,\mathrm{m}$ after twenty days of regular growth, how tall was it at the beginning of the first day?}
{$4\,\mathrm{m}$}{$5.3\,\mathrm{m}$}{$2.7\,\mathrm{m}$}{$10\,\mathrm{m}$}{$4.3\,\mathrm{m}$}{}}
\LANGpl{
\Question{
Pewien gatunek bambusa rośnie  $1{,}3\,\mathrm{m}$ dziennie podczas okresu wegetacji. Wiedząc, że  osiągnął wysokość $30\,\mathrm{m}$ apo dwudziestu dniach regularnego wzrostu oblicz jak wysoki był na początku pierwszego dnia.}
{$4\,\mathrm{m}$}{$5{,}3\,\mathrm{m}$}{$2{,}7\,\mathrm{m}$}{$10\,\mathrm{m}$}{$4{,}3\,\mathrm{m}$}{}}
\LANGcs{
\Question{
Jistý druh bambusu vyroste ve vegetačním období za den o $1{,}3\,\mathrm{m}$. Kolik měřil na začátku prvního dne měření, když po dvaceti dnech pravidelného růstu dosáhl výšky $30\,\mathrm{m}$?}
{$4\,\mathrm{m}$}{$5{,}3\,\mathrm{m}$}{$2{,}7\,\mathrm{m}$}{$10\,\mathrm{m}$}{$4{,}3\,\mathrm{m}$}{}}
\LANGsk{
\Question{
Istý druh bambusu vyrastie vo vegetačnom období za deň o $1{,}3\,\mathrm{m}$. Koľko meral na začiatku prvého dňa merania, keď po dvadsiatich dňoch pravidelného rastu dosiahol výšku $30\,\mathrm{m}$?}
{$4\,\mathrm{m}$}{$5{,}3\,\mathrm{m}$}{$2{,}7\,\mathrm{m}$}{$10\,\mathrm{m}$}{$4{,}3\,\mathrm{m}$}{}}
\LANGes{
\Question{
Un cierto tipo de bambú crece $1.3\,\mathrm{m}$ por día durante el período de vegetación. ¿Cuántos metros midía el primer día si sabemos que después de veinte días de crecimiento regular alcanzó una altura de $30\,\mathrm{m}$?}
{$4\,\mathrm{m}$}{$5.3\,\mathrm{m}$}{$2.7\,\mathrm{m}$}{$10\,\mathrm{m}$}{$4.3\,\mathrm{m}$}{}}
\End

\Data\ID{1003047802}
\Updated{2021-05-03 08:05:49}
\Level{C}
\SubArea{Arithmetic sequences}
\LANGen{
\Question{
A student solves several math problems every day. How many problems will he solve in $14$ days if he manages to solve $5$ on the first day and then gradually increases the number of problems solved by $2$ a day?}
{$252$}{$140$}{$434$}{$364$}{$70$}{}}
\LANGpl{
\Question{
Student rozwiązuje kilka zadań matematycznych dziennie. Ile zadań rozwiąże w  $14$ dni jeśli zdołał rozwiązać $5$ pierwszego dnia, a każdego kolejnego dnia ta liczba wzrasta o $2$?}
{$252$}{$140$}{$434$}{$364$}{$70$}{}}
\LANGcs{
\Question{
Student vypočítá každý den několik příkladů. Kolik by jich vypočítal za $14$ dní, kdyby první den vypočítal $5$ a pak postupně zvyšoval počet vypočítaných příkladů každý den o $2$ ?}
{$252$}{$140$}{$434$}{$364$}{$70$}{}}
\LANGsk{
\Question{
Študent vypočíta každý deň niekoľko príkladov. Koľko by ich vypočítal za $14$ dní, keby prvý deň vypočítal $5$ a potom postupne zvyšoval počet vypočítaných príkladov každý deň o $2$ ?}
{$252$}{$140$}{$434$}{$364$}{$70$}{}}
\LANGes{
\Question{
Un estudiante resuelve varios ejercicios de matemáticas todos los días. ¿Cuántos ejercicios resolvería en $14$ días si fuera capaz de resolver $5$ ejercicios el primer día y luego aumentara gradualmente el número de ejercicios resueltos en $2$ al día?}
{$252$}{$140$}{$434$}{$364$}{$70$}{}}
\End

\Data\ID{1003047803}
\Updated{2021-05-03 08:05:55}
\Level{C}
\SubArea{Arithmetic sequences}
\LANGen{
\Question{
In Duolingo application, every user earns so-called lingots (=virtual currency), providing he or she learns at least $10$ minutes for $10$ consecutive days. For the first $10$ days a user  earns $1$ lingot, $2$ lingots for the next ten-day period, $3$ lingots for the following $10$ days (i.e. for the first $30$ days a user  earns $6$ lingots), etc. Find the smallest number of days in which a user can accumulate $1000$ lingots.}
{$450$}{$45$}{$440$}{$44$}{$430$}{}}
\LANGpl{
\Question{
W aplikacji Duolingo każdy użytkownik zarabia tzw. lingoty (=wirtualną walutę), zakładając, że uczy się po  $10$ minut przez $10$ kolejnych dni. Przez pierwsze $10$ dni użytkownik zarabia $1$ lingot, $2$ lingoty przez kolejny dziesięciodniowy okres, $3$ lingoty przez kolejne $10$ dni (tzn. przez pierwsze $30$ użytkownik zyskuje  $6$ lingoty ), itd. Wyznacz najmniejszą liczbę dni, w ciągu których użytkownik może zebrać $1000$ lingotów.}
{$450$}{$45$}{$440$}{$44$}{$430$}{}}
\LANGcs{
\Question{
V aplikaci Duolingo získává uživatel za aspoň $10$ minut učení po dobu $10$ dní bez přerušení lingoty (=virtuální měna). Za prvních $10$ dní $1$ lingot, za druhých $10$ dní $2$ lingoty, za dalších $10$ dní $3$ lingoty (tj. za prvních $30$ dní získá $6$ lingotů), atd. Za kolik nejméně dní může uživatel nashromáždit $1000$ lingotů?}
{$450$}{$45$}{$440$}{$44$}{$430$}{}}
\LANGsk{
\Question{
V aplikácii Duolingo získava užívateľ za aspoň $10$ minút učenia po dobu $10$ dní bez prerušenia lingoty (=virtuálna mena). Za prvých $10$ dní $1$ lingot, za druhých $10$ dní $2$ lingoty, za ďalších $10$ dní $3$ lingoty (tj. za prvých $30$ dní získa $6$ lingot), atď. Za koľko najmenej dní môže užívateľ nazhromaždiť $1000$ lingot?}
{$450$}{$45$}{$440$}{$44$}{$430$}{}}
\LANGes{
\Question{
En la aplicación Duolingo, cada usuario gana los llamados lingots (=moneda virtual), siempre que estudie al menos $10$ minutos durante $10$ días consecutivos. Durante los $10$ primeros días, un usuario gana $1$ lingot, $2$ lingots por el próximo período de diez días, $3$ lingots por los siguientes $10$ días (es decir, por los $30$ primeros días, un usuario gana $6$ lingots), etc. Halla el menor número de días que tiene que estudiar un usuario para ganar  $1000$ lingots.}
{$450$}{$45$}{$440$}{$44$}{$430$}{}}
\End

\Data\ID{1003047804}
\Updated{2021-05-03 08:05:08}
\Level{C}
\SubArea{Arithmetic sequences}
\LANGen{
\Question{
The pianist wanted to learn a new composition in $3$ weeks ($21$ days). He decided to learn  the same number of bars (measures) a day. In the end, however, he met his plan only on the first day. Every other day, he managed to learn by one bar fewer than the previous day. Find how many bars did he learn on the $15$th day knowing that he was able to learn a total of $462$ measures (in $21$ days).}
{$18$}{$22$}{$32$}{$15$}{$20$}{}}
\LANGpl{
\Question{
Pianista chciał nauczyć się nowej kompozycji w ciągu  $3$ tygodni ($21$ dni). Zdecydował, że będzie uczył się tej samej liczby taktów każdego dnia. Ostatecznie udało mu się wypełnić plan jedynie pierwszego dnia. Każdego kolejnego dnia udało mu się nauczyć o jedne takt mniej niż każdego poprzedniego dnia. Ilu taktów nauczył się  $15$ dnia wiedząc, że całościowo zdołał opanować $462$ taktów (w $21$ dni).}
{$18$}{$22$}{$32$}{$15$}{$20$}{}}
\LANGcs{
\Question{
Klavírista chtěl nastudovat novou skladbu za $3$ týdny ($21$ dní). Rozhodl se každý den nacvičit stejný počet taktů. Nakonec ale plán dodržel jen první den. Každý další den nacvičil o jeden takt méně než předchozí den. Kolik taktů nacvičil $15$. den, když (za  $21$ dní) stihl nacvičit celkem $462$ taktů?}
{$18$}{$22$}{$32$}{$15$}{$20$}{}}
\LANGsk{
\Question{
Klavirista chcel naštudovať novú skladbu za $3$ týždne ($21$ dní). Rozhodol sa každý deň nacvičiť rovnaký počet taktov. Nakoniec ale plán dodržal len prvý deň. Každý ďalší deň nacvičil o jeden takt menej než predchádzajúci deň. Koľko taktov nacvičil $15$. deň, keď (za  $21$ dní) stihol nacvičiť celkom $462$ taktov?}
{$18$}{$22$}{$32$}{$15$}{$20$}{}}
\LANGes{
\Question{
Un pianista quería aprender una nueva composición en $3$ semanas ($21$ días). Decidió aprender la misma cantidad de compases al día. Al final, sin embargo, cumplió su plan solo el primer día. Después, cada día, logró aprender un compás menos que el día anterior. Halla cuántos compases aprendió el día $15$ si sabemos que aprendió $462$ compases en $21$ días.}
{$18$}{$22$}{$32$}{$15$}{$20$}{}}
\End

\Data\ID{1003047805}
\Updated{2021-05-03 08:05:23}
\Level{C}
\SubArea{Arithmetic sequences}
\LANGen{
\Question{
A cyclist plans to travel $1666\,\mathrm{km}$ in $14$ days during her vacation. She knows that the number of kilometers she rides every day will decrease by the same number so she planned her route accordingly. At the beginning  of the last day she was only $80\,\mathrm{km}$ away from her goal. What is the difference between the numbers of kilometers she rode on two consecutive days?}
{$6\,\mathrm{km}$}{$7\,\mathrm{km}$}{$5\,\mathrm{km}$}{$4\,\mathrm{km}$}{$3\,\mathrm{km}$}{}}
\LANGpl{
\Question{
Rowerzysta planuje przejechać $1666\,\mathrm{km}$ w $14$ dni podczas wakacji. Wie, że liczba kilometrów przebytych każdego dnia będzie zmniejszać się o tę samą liczbę, zatem zaplanowała trasę zgodnie z posiadaną wiedzą. Na początku ostatniego dnia była zaledwie  $80\,\mathrm{km}$ od celu podróży. Jaka jest różnica pomiędzy liczbą kilometrów przebytych w ciągu kolejnych dwóch dni?}
{$6\,\mathrm{km}$}{$7\,\mathrm{km}$}{$5\,\mathrm{km}$}{$4\,\mathrm{km}$}{$3\,\mathrm{km}$}{}}
\LANGcs{
\Question{
Cyklista má v plánu ujet $1666\,\mathrm{km}$ za $14$ dní dovolené. Ví, že postupně ujede každý den o stejný počet kilometrů méně než předchozí, a podle toho si naplánoval trasu. Poslední den mu zbývalo ujet jen $80\,\mathrm{km}$. Jaký je rozdíl v ujetých kilometrech mezi dvěma po sobě jdoucími dny?}
{$6\,\mathrm{km}$}{$7\,\mathrm{km}$}{$5\,\mathrm{km}$}{$4\,\mathrm{km}$}{$3\,\mathrm{km}$}{}}
\LANGsk{
\Question{
Cyklista má v pláne prejsť $1666\,\mathrm{km}$ za $14$ dní dovolenky. Vie, že postupne prejde každý deň o rovnaký počet kilometrov menej než predchádzajúci, a podľa toho si naplánoval trasu. Posledný deň mu zostalo prejsť len $80\,\mathrm{km}$. Aký je rozdiel v prejdených kilometroch medzi dvoma po sebe idúcimi dňami?}
{$6\,\mathrm{km}$}{$7\,\mathrm{km}$}{$5\,\mathrm{km}$}{$4\,\mathrm{km}$}{$3\,\mathrm{km}$}{}}
\LANGes{
\Question{
Una ciclista planea viajar $1666\,\mathrm{km}$ en $14$ días durante sus vacaciones. Ella sabe que la cantidad de kilómetros que recorre todos los días disminuirá la misma cantidad que el día anterior, y a base de eso planeó su ruta. Al comienzo del último día, estaba a solo $80\,\mathrm{km}$ de distancia de su objetivo. ¿Cuál es la diferencia entre la cantidad de kilómetros recorridos en dos días consecutivos?}
{$6\,\mathrm{km}$}{$7\,\mathrm{km}$}{$5\,\mathrm{km}$}{$4\,\mathrm{km}$}{$3\,\mathrm{km}$}{}}
\End

\Data\ID{1003047806}
\Updated{2021-05-03 08:05:38}
\Level{C}
\SubArea{Arithmetic sequences}
\LANGen{
\Question{
Every school must pay a registration fee for each participant it sends to a math competition. The fee for the first participant is $10$ euros, for each additional one the fee is one euro less. The maximum number of participants any school can register is $10$. Find the relationship between the price ($c$) paid by the school  and the number of students registered ($n$).}
{$c=\frac n2(21-n)$}{$c=10-\frac{n^2}2$}{$c=\frac{11n}2$}{$c=\frac n2(10+10n)$}{$c=\frac n2(11-n)$}{}}
\LANGpl{
\Question{
Każda szkoła musi wpłacić wpisowe za każdego uczestnika konkursu matematycznego. Opłata za pierwszego uczestnika wynosi $10$ euro, dla każdego kolejnego opłata jest o euro niższa. Maksymalna liczba uczestników, którą może zglosić szkoła wynosi $10$. Jaka jest zależność pomiędzy opłatą ($c$) wniesioną przez szkołę a liczbą zapisanych studentów ($n$).}
{$c=\frac n2(21-n)$}{$c=10-\frac{n^2}2$}{$c=\frac{11n}2$}{$c=\frac n2(10+10n)$}{$c=\frac n2(11-n)$}{}}
\LANGcs{
\Question{
Při přihlašování studentů na matematickou soutěž platí škola za každého účastníka registrační poplatek. Za prvního přihlášeného platí $10$ euro, za každého dalšího o euro méně, více než $10$ studentů nesmí škola přihlásit. Vyjádřete vztah závislosti ceny ($c$), kterou škola zaplatí, na počtu ($n$) přihlášených studentů.}
{$c=\frac n2(21-n)$}{$c=10-\frac{n^2}2$}{$c=\frac{11n}2$}{$c=\frac n2(10+10n)$}{$c=\frac n2(11-n)$}{}}
\LANGsk{
\Question{
Pri prihlasovaní študentov na matematickú súťaž platí škola za každého účastníka registračný poplatok. Za prvého prihláseného platí $10$ euro, za každého ďalšieho o euro menej, viac než $10$ študentov nesmie škola prihlásiť. Vyjadrite vzťah závislosti ceny ($c$), ktorú škola zaplatí, na počtu ($n$) prihlásených študentov.}
{$c=\frac n2(21-n)$}{$c=10-\frac{n^2}2$}{$c=\frac{11n}2$}{$c=\frac n2(10+10n)$}{$c=\frac n2(11-n)$}{}}
\LANGes{
\Question{
Cada escuela debe pagar una tarifa de inscripción por cada participante enviado a un concurso de matemáticas. La tarifa para el primer participante es de $10$ euros, por cada uno más, la tarifa es un euro menos. El número máximo de participantes de una escuela es $10$. Halla la relación entre el precio ($c$) pagado por la escuela y el número de estudiantes inscritos ($n$).}
{$c=\frac n2(21-n)$}{$c=10-\frac{n^2}2$}{$c=\frac{11n}2$}{$c=\frac n2(10+10n)$}{$c=\frac n2(11-n)$}{}}
\End

\Data\ID{1003047807}
\Updated{2020-07-15 03:07:12}
\Level{C}
\SubArea{Arithmetic sequences}
\LANGen{
\Question{
At a dorm, students play with rolls of toilet paper, building a pyramid of it. There is one roll at the top of the pyramid, in each lower row there is one roll more than in the one above it. How high will the pyramid be, given the fact that they have a total of $171$ rolls and one roll is $9.5\,\mathrm{cm}$ high?}
{$171\,\mathrm{cm}$}{$2\,\mathrm{m}$}{$180\,\mathrm{cm}$}{$95\,\mathrm{cm}$}{$123.5\,\mathrm{cm}$}{}}
\LANGpl{
\Question{
W internacie studenci bawią się rolkami papieru toaletowego budując z nich piramidę. Na szczycie piramidy znajduje się jedna rolka, a w każdym kolejnym rzędzie znajduje się o jedną rolkę więcej niż w rzędzie poprzednim. Jak wysoka będzie piramida, jeśli dysponują  $171$ rolkami, a każda z nich ma wysokość  $9{,}5\,\mathrm{cm}$?}
{$171\,\mathrm{cm}$}{$2\,\mathrm{m}$}{$180\,\mathrm{cm}$}{$95\,\mathrm{cm}$}{$123{,}5\,\mathrm{cm}$}{}}
\LANGcs{
\Question{
Na kolejích si studenti staví pyramidu z ruliček od toaletního papíru. Nahoře je jedna rulička, v každé další řadě je o ruličku více. Jak bude pyramida vysoká, jestliže mají celkem $171$ ruliček a jedna rulička měří $9{,}5\,\mathrm{cm}$?}
{$171\,\mathrm{cm}$}{$2\,\mathrm{m}$}{$180\,\mathrm{cm}$}{$95\,\mathrm{cm}$}{$123{,}5\,\mathrm{cm}$}{}}
\LANGsk{
\Question{
Na internáte si študenti stavajú pyramídu z roliek od toaletného papiera. Navrchu je jedna rolka, v každom ďalšom rade je o rolku viac. Aká bude pyramída vysoká, ak majú celkom $171$ roliek a jedna rolka meria $9{,}5\,\mathrm{cm}$?}
{$171\,\mathrm{cm}$}{$2\,\mathrm{m}$}{$180\,\mathrm{cm}$}{$95\,\mathrm{cm}$}{$123{,}5\,\mathrm{cm}$}{}}
\LANGes{
\Question{
En un dormitorio, los estudiantes juegan con rollos de papel higiénico construyendo una pirámide. En la cúspide de la pirámide hay un rollo, en cada fila inferior hay un rollo más que en la de arriba. ¿Cuánto mide la pirámide suponiendo que consta de $171$ rollos y un rollo mide $9.5\,\mathrm{cm}$?}
{$171\,\mathrm{cm}$}{$2\,\mathrm{m}$}{$180\,\mathrm{cm}$}{$95\,\mathrm{cm}$}{$123.5\,\mathrm{cm}$}{}}
\End

\Data\ID{1003047808}
\Updated{2020-07-15 03:07:12}
\Level{C}
\SubArea{Arithmetic sequences}
\LANGen{
\Question{
A heavy smoker wants to reduce his daily consumption by $2$ cigarettes every $30$ days. How much will he save in $360$ days given the fact that one box of cigarettes ($20$ cigarettes) costs $80\ \mathrm{CZK}$?}
{$18\,720\ \mathrm{CZK}$}{$624\ \mathrm{CZK}$}{$4\,680\ \mathrm{CZK}$}{$12\,480\ \mathrm{CZK}$}{$6\,240\ \mathrm{CZK}$}{}}
\LANGpl{
\Question{
Nałogowy palacz chce zmniejszyć ilość wypalanych papierosów o  $2$ co $30$ dni. Ile zaoszczędzi w ciągu  $360$ dni jeśli jedna paczka papierosów  ($20$ papierosów) kosztuje $80\ \mathrm{CZK}$?}
{$18\,720\ \mathrm{CZK}$}{$624\ \mathrm{CZK}$}{$4\,680\ \mathrm{CZK}$}{$12\,480\ \mathrm{CZK}$}{$6\,240\ \mathrm{CZK}$}{}}
\LANGcs{
\Question{
Náruživý kuřák chce každých $30$ dní snížit svoji denní spotřebu o $2$ cigarety. Kolik ušetří za $360$ dní, jestliže jedna krabička  ($20$ cigaret) stojí $80\ \mathrm{CZK}$?}
{$18\,720\ \mathrm{CZK}$}{$624\ \mathrm{CZK}$}{$4\,680\ \mathrm{CZK}$}{$12\,480\ \mathrm{CZK}$}{$6\,240\ \mathrm{CZK}$}{}}
\LANGsk{
\Question{
Náruživý fajčiar chce každých $30$ dní znížiť svoju denní spotrebu o $2$ cigarety. Koľko ušetrí za $360$ dní, ak jedna krabička  ($20$ cigariet) stojí $80\ \mathrm{CZK}$?}
{$18\,720\ \mathrm{CZK}$}{$624\ \mathrm{CZK}$}{$4\,680\ \mathrm{CZK}$}{$12\,480\ \mathrm{CZK}$}{$6\,240\ \mathrm{CZK}$}{}}
\LANGes{
\Question{
Un fumador apasionado quiere reducir su consumo diario en $2$ cigarrillos cada $30$ días. ¿Cuánto dinero ahorrará en $360$ días suponiendo que una cajetilla ($20$ cigarrillos) cuesta $80\ \mathrm{CZK}$?}
{$18\,720\ \mathrm{CZK}$}{$624\ \mathrm{CZK}$}{$4\,680\ \mathrm{CZK}$}{$12\,480\ \mathrm{CZK}$}{$6\,240\ \mathrm{CZK}$}{}}
\End

\Data\ID{1003107201}
\Updated{2020-07-15 03:07:12}
\Level{C}
\SubArea{Arithmetic sequences}
\LANGen{
\Question{
The sum of the first ten terms of an arithmetic sequence with odd subscripts is $190$, the sum of the first ten terms with even subscripts is $230$. Find the first term.}
{$-17$}{$-34$}{$5$}{$4$}{$10$}{}}
\LANGpl{
\Question{
Suma pierwszych dziesięciu wyrazów ciągu arytmetycznego z nieparzystymi z indeksem dolnym wynosi  $190$, suma pierwszych dziesięciu wyrazów  z parzystymi indeksami dolnymi wynosi $230$. Wyznacz pierwszy wyraz.}
{$-17$}{$-34$}{$5$}{$4$}{$10$}{}}
\LANGcs{
\Question{
Součet prvních deseti členů aritmetické posloupnosti s lichými indexy je $190$, součet prvních deseti členů se sudými indexy je $230$. Určete první člen.}
{$-17$}{$-34$}{$5$}{$4$}{$10$}{}}
\LANGsk{
\Question{
Súčet prvých desiatich členov aritmetickej postupnosti s nepárnymi indexmi je $190$, súčet prvých desiatich členov s párnymi indexmi je $230$. Určte prvý člen.}
{$-17$}{$-34$}{$5$}{$4$}{$10$}{}}
\LANGes{
\Question{
La suma de los diez primeros términos de una progresión aritmética con subíndices impares es $190$, la suma de los diez primeros términos con subíndices pares es $230$. Halla el primer término.}
{$-17$}{$-34$}{$5$}{$4$}{$10$}{}}
\End

\Data\ID{1003107202}
\Updated{2021-05-03 08:05:21}
\Level{C}
\SubArea{Arithmetic sequences}
\LANGen{
\Question{
The second term of an arithmetic sequence is $-21$, the fifth term is $21$. How many of the first consecutive terms do we have to add up to get a sum of more than $50$?}
{$8$}{$7$}{$6$}{$9$}{$5$}{}}
\LANGpl{
\Question{
Drugi wyraz ciągu arytmetycznego wynosi  $-21$, piąty jest równy $21$. Ile spośród kolejnych wyrazów tego ciągu musimy dodać, by otrzymać sumę wyższą niż  $50$?}
{$8$}{$7$}{$6$}{$9$}{$5$}{}}
\LANGcs{
\Question{
Druhý člen aritmetické posloupnosti je $-21$, pátý člen je $21$. Kolik prvních členů musíme minimálně sečíst, aby byl součet větší než $50$?}
{$8$}{$7$}{$6$}{$9$}{$5$}{}}
\LANGsk{
\Question{
Druhý člen aritmetickej postupnosti je $-21$, piaty člen je $21$. Koľko prvých členov musíme minimálne sčítať, aby bol súčet väčší ako $50$?}
{$8$}{$7$}{$6$}{$9$}{$5$}{}}
\LANGes{
\Question{
El segundo término de una progresión aritmética es $-21$, el quinto término es $21$. ¿Empezando por el principio, cuántos términos consecutivos tenemos que sumar para que la suma sea mayor que $50$?}
{$8$}{$7$}{$6$}{$9$}{$5$}{}}
\End

\Data\ID{1003107203}
\Updated{2021-05-03 08:05:02}
\Level{C}
\SubArea{Arithmetic sequences}
\LANGen{
\Question{
Insert such $3$ numbers between the roots of the equation $x^2-10x-119=0$, so that together with the roots of the equation they form $5$ consecutive arithmetic sequence terms. What is the central term?}
{$5$}{$17$}{$6$}{$-1$}{$-7$}{}}
\LANGpl{
\Question{
Wstaw takie  $3$ liczy pomiędzy pierwiastki równania $x^2-10x-119=0$, tak by razem z pierwiastkami równania utworzyły  $5$ kolejnych wyrazów ciągu arytmetycznego. Jaki jest środkowy wyraz?}
{$5$}{$17$}{$6$}{$-1$}{$-7$}{}}
\LANGcs{
\Question{
Mezi kořeny rovnice $x^2-10x-119=0$ vložte $3$ čísla tak, aby spolu s kořeny rovnice tvořila $5$ následujících členů aritmetické posloupnosti. Prostřední z nich je rovno:}
{$5$}{$17$}{$6$}{$-1$}{$-7$}{}}
\LANGsk{
\Question{
Medzi korene rovnice $x^2-10x-119=0$ vložte $3$ čísla tak, aby spolu s koreňmi rovnice tvorili $5$ nasledujúcich členov aritmetickej postupnosti. Prostredný z nich je rovný:}
{$5$}{$17$}{$6$}{$-1$}{$-7$}{}}
\LANGes{
\Question{
Entre las raíces de la siguiente ecuación $x^2-10x-119=0$ introduce $3$ números para que junto con las raíces de la ecuación formen $5$ términos consecutivos de una prograsión aritmética. ¿Cuál es el número central?}
{$5$}{$17$}{$6$}{$-1$}{$-7$}{}}
\End

\Data\ID{1003107204}
\Updated{2021-05-03 08:05:05}
\Level{C}
\SubArea{Arithmetic sequences}
\LANGen{
\Question{
Positive numbers which when divided by $3$ leave a remainder of $2$ are terms of an increasing arithmetic sequence. The sum of the first $15$ terms is $480$. Find the $3$rd term of this sequence.}
{$17$}{$11$}{$20$}{$5$}{$23$}{}}
\LANGpl{
\Question{
Liczby dodatnie podzielone przez  $3$ z resztą  $2$ są wyrazami rosnącego ciągu arytmetycznego. Suma pierwszych $15$ wyrazów wynosi $480$. Wyznacz  $3$ wyraz tego ciągu.}
{$17$}{$11$}{$20$}{$5$}{$23$}{}}
\LANGcs{
\Question{
Členy rostoucí aritmetické posloupnosti s kladnými členy jsou čísla, která při dělení $3$ dávají zbytek $2$. Určete třetí člen, jestliže součet prvních $15$ členů je $480$.}
{$17$}{$11$}{$20$}{$5$}{$23$}{}}
\LANGsk{
\Question{
Členy rastúcej aritmetickej postupnosti s kladnými členmi sú čísla, ktoré pri delení $3$ dávajú zvyšok $2$. Určte tretí člen, ak súčet prvých $15$ členov je $480$.}
{$17$}{$11$}{$20$}{$5$}{$23$}{}}
\LANGes{
\Question{
Los números positivos que al dividirse entre $3$ dejan un resto de $2$ son términos de una progresión aritmética creciente. La suma de los $15$ primeros términos es $480$. Halla el tercer término de la progresión.}
{$17$}{$11$}{$20$}{$5$}{$23$}{}}
\End

\Data\ID{1003107205}
\Updated{2021-05-03 08:05:08}
\Level{C}
\SubArea{Arithmetic sequences}
\LANGen{
\Question{
Sizes of angles in a triangle form three consecutive arithmetic sequence terms. The size of the largest angle is four times the smallest one. Find the size of the smallest angle of the triangle.}
{$24^{\circ}$}{$30^{\circ}$}{$60^{\circ}$}{$20^{\circ}$}{$35^{\circ}$}{}}
\LANGpl{
\Question{
Miary kątów w trójkącie tworzą trzy kolejne wyrazy ciągu arytmetycznego. Miara największego jest równa cztery razy miara najmniejszego z nich  Jaka jest miara najmniejszego kąta?}
{$24^{\circ}$}{$30^{\circ}$}{$60^{\circ}$}{$20^{\circ}$}{$35^{\circ}$}{}}
\LANGcs{
\Question{
Velikosti úhlů v trojúhelníku tvoří tři po sobě jdoucí členy aritmetické posloupnosti. Velikost největšího z nich je čtyřnásobek velikosti nejmenšího. Určete velikost nejmenšího úhlu trojúhelníku.}
{$24^{\circ}$}{$30^{\circ}$}{$60^{\circ}$}{$20^{\circ}$}{$35^{\circ}$}{}}
\LANGsk{
\Question{
Veľkosť uhlov v trojuholníku tvoria tri po sebe idúce členy aritmetickej postupnosti. Veľkosť najväčšieho z nich je štvornásobok veľkosti najmenšieho. Určte veľkosti najmenšieho uhla v trojuholníku.}
{$24^{\circ}$}{$30^{\circ}$}{$60^{\circ}$}{$20^{\circ}$}{$35^{\circ}$}{}}
\LANGes{
\Question{
Las medidas de los ángulos de un triángulo forman tres términos consecutivos de una progresión aritmética. La medida del ángulo más grande es el cuádruple de la del más pequeño. Calcula la medida del ángulo más pequeño del triángulo.}
{$24^{\circ}$}{$30^{\circ}$}{$60^{\circ}$}{$20^{\circ}$}{$35^{\circ}$}{}}
\End

\Data\ID{1003107206}
\Updated{2021-05-03 08:05:30}
\Level{C}
\SubArea{Arithmetic sequences}
\LANGen{
\Question{
Solve the equation with the natural variable $n$: 
$$ 5+8+11+14+\dots+(3n+2)=735 $$}
{$21$}{$65$}{$20$}{$68$}{$10$}{}}
\LANGpl{
\Question{
Rozwiąż równanie z naturalną zmienną $n$: 
$$ 5+8+11+14+\dots+(3n+2)=735 $$}
{$21$}{$65$}{$20$}{$68$}{$10$}{}}
\LANGcs{
\Question{
Řešte rovnici s neznámou $n\in\mathbb{N}$: 
$$ 5+8+11+14+\dots+(3n+2)=735 $$}
{$21$}{$65$}{$20$}{$68$}{$10$}{}}
\LANGsk{
\Question{
Riešte rovnicu s neznámou $n\in\mathbb{N}$: 
$$ 5+8+11+14+\dots+(3n+2)=735 $$}
{$21$}{$65$}{$20$}{$68$}{$10$}{}}
\LANGes{
\Question{
Resuelve la ecuación con variable natural $n$: 
$$ 5+8+11+14+\dots+(3n+2)=735 $$}
{$21$}{$65$}{$20$}{$68$}{$10$}{}}
\End

\Data\ID{1003107207}
\Updated{2020-07-15 03:07:12}
\Level{C}
\SubArea{Arithmetic sequences}
\LANGen{
\Question{
Find the set of all solutions of the given inequality with the natural variable $x$:
$$ 5x+10x+15x+\dots+50x \leq 1000 $$}
{$ \{1;2;3\} $}{$ \{ 2 \} $}{$\{ 1;2;3;4;5\} $}{$\{ 1 \}$}{$\{1;2\}$}{}}
\LANGpl{
\Question{
Wyznacz zbiór wszystkich rozwiązań podanej nierówności z naturalna zmienną  $x$:
$$ 5x+10x+15x+\dots+50x \leq 1000 $$}
{$ \{1;2;3\} $}{$ \{ 2 \} $}{$\{ 1;2;3;4;5\} $}{$\{ 1 \}$}{$\{1;2\}$}{}}
\LANGcs{
\Question{
Najděte množinu všech řešení nerovnice s neznámou $x\in\mathbb{N}$:
$$ 5x+10x+15x+\dots+50x \leq 1000 $$}
{$ \{1;2;3\} $}{$ \{ 2 \} $}{$\{ 1;2;3;4;5\} $}{$\{ 1 \}$}{$\{1;2\}$}{}}
\LANGsk{
\Question{
Nájdite množinu všetkých riešení nerovnice s neznámou $x\in\mathbb{N}$:
$$ 5x+10x+15x+\dots+50x \leq 1000 $$}
{$ \{1;2;3\} $}{$ \{ 2 \} $}{$\{ 1;2;3;4;5\} $}{$\{ 1 \}$}{$\{1;2\}$}{}}
\LANGes{
\Question{
Halla el conjunto de todas las soluciones de la siguiente inecuación con la variable natural $x$:
$$ 5x+10x+15x+\dots+50x \leq 1000 $$}
{$ \{1;2;3\} $}{$ \{ 2 \} $}{$\{ 1;2;3;4;5\} $}{$\{ 1 \}$}{$\{1;2\}$}{}}
\End

\Data\ID{1003107208}
\Updated{2020-07-15 03:07:12}
\Level{C}
\SubArea{Arithmetic sequences}
\LANGen{
\Question{
Solve the equation with the natural  variable $x$:
$$ \sum\limits_{n=1}^x(-3n+10)=-45 $$}
{$9$}{$5$}{$10$}{$13$}{$12$}{}}
\LANGpl{
\Question{
Rozwiąż równanie z naturalną zmienną $x$:
$$ \sum\limits_{n=1}^x(-3n+10)=-45 $$}
{$9$}{$5$}{$10$}{$13$}{$12$}{}}
\LANGcs{
\Question{
Řešte rovnici s neznámou $\ x\in\mathbb{N}$:
$$ \sum\limits_{n=1}^x(-3n+10)=-45 $$}
{$9$}{$5$}{$10$}{$13$}{$12$}{}}
\LANGsk{
\Question{
Riešte rovnicu s neznámou $\ x\in\mathbb{N}$:
$$ \sum\limits_{n=1}^x(-3n+10)=-45 $$}
{$9$}{$5$}{$10$}{$13$}{$12$}{}}
\LANGes{
\Question{
Resuelve la ecuación con la variable natural $x$:
$$ \sum\limits_{n=1}^x(-3n+10)=-45 $$}
{$9$}{$5$}{$10$}{$13$}{$12$}{}}
\End

\Data\ID{1003107209}
\Updated{2020-07-15 03:07:12}
\Level{C}
\SubArea{Arithmetic sequences}
\LANGen{
\Question{
Solve the equation with the real variable $x$:
$$ \sum\limits_{n=1}^{23} xn=92 $$}
{$\frac13$}{$4$}{$\frac14$}{$1$}{$\frac12$}{}}
\LANGpl{
\Question{
Rozwiąż równanie z rzeczywistą zmienną$x$:
$$ \sum\limits_{n=1}^{23} xn=92 $$}
{$\frac13$}{$4$}{$\frac14$}{$1$}{$\frac12$}{}}
\LANGcs{
\Question{
Řešte rovnici s neznámou $\ x\in\mathbb{R}$:
$$ \sum\limits_{n=1}^{23} xn=92 $$}
{$\frac13$}{$4$}{$\frac14$}{$1$}{$\frac12$}{}}
\LANGsk{
\Question{
Riešte rovnicu s neznámou $\ x\in\mathbb{R}$:
$$ \sum\limits_{n=1}^{23} xn=92 $$}
{$\frac13$}{$4$}{$\frac14$}{$1$}{$\frac12$}{}}
\LANGes{
\Question{
Resuelve la ecuación con la variable natural $x$:
$$ \sum\limits_{n=1}^{23} xn=92 $$}
{$\frac13$}{$4$}{$\frac14$}{$1$}{$\frac12$}{}}
\End

\Data\ID{1003107210}
\Updated{2020-12-28 03:12:30}
\Level{C}
\SubArea{Arithmetic sequences}
\LANGen{
\Question{
Solve the inequality with the natural variable $x$:
$$\sum\limits_{n=1}^x \left(\frac25n-1\right) \leq x $$}
{$ x\leq9;\ x\in\mathbb{N}$}{$x\leq10;\ x\in\mathbb{N}$}{$x\in(-\infty;8]$}{$x\geq8;\ x\in\mathbb{N}$}{there is no solution in $\mathbb{N}$}{}}
\LANGpl{
\Question{
Rozwiąż nierówność z naturalną zmienną $x$:
$$\sum\limits_{n=1}^x \left(\frac25n-1\right) \leq x $$}
{$ x\leq9;\ x\in\mathbb{N}$}{$x\leq10;\ x\in\mathbb{N}$}{$x\in(-\infty;8\rangle$}{$x\geq8;\ x\in\mathbb{N}$}{nie ma rozwiązania w  $\mathbb{N}$}{}}
\LANGcs{
\Question{
Řešte nerovnici s neznámou $\ x\in\mathbb{N}$:
$$\sum\limits_{n=1}^x \left(\frac25n-1\right) \leq x $$}
{$ x\leq9;\ x\in\mathbb{N}$}{$x\leq10;\ x\in\mathbb{N}$}{$x\in(-\infty;8\rangle$}{$x\geq8;\ x\in\mathbb{N}$}{nemá řešení v $\mathbb{N}$}{}}
\LANGsk{
\Question{
Riešte nerovnicu s neznámou $\ x\in\mathbb{N}$:
$$\sum\limits_{n=1}^x \left(\frac25n-1\right) \leq x $$}
{$ x\leq9;\ x\in\mathbb{N}$}{$x\leq10;\ x\in\mathbb{N}$}{$x\in(-\infty;8\rangle$}{$x\geq8;\ x\in\mathbb{N}$}{nemá riešenie v $\mathbb{N}$}{}}
\LANGes{
\Question{
Resuelve la ecuación con la variable natural $x$:
$$\sum\limits_{n=1}^x \left(\frac25n-1\right) \leq x $$}
{$ x\leq9;\ x\in\mathbb{N}$}{$x\leq10;\ x\in\mathbb{N}$}{$x\in(-\infty;8]$}{$x\geq8;\ x\in\mathbb{N}$}{there is no solution in $\mathbb{N}$}{}}
\End

\Data\ID{9000064801}
\Updated{2020-07-15 03:07:37}
\Level{C}
\SubArea{Arithmetic sequences}
\LANGen{
\Question{
The sides of a right triangle form three consecutive terms of an arithmetic sequence. The perimeter
of the triangle is \(60\, \mathrm{cm}\).
Find the hypotenuse of the triangle (the longest side).}
{\(25\, \mathrm{cm}\)}{\(12\, \mathrm{cm}\)}{\(15\, \mathrm{cm}\)}{\(20\, \mathrm{cm}\)}{\(30\, \mathrm{cm}\)}{}}
\LANGpl{
\Question{
Trzy kolejne wyrazy ciągu arytmetycznego są długościami boków trójkąta prostokątnego. Obwód tego trójkąta wynosi   \(60\, \mathrm{cm}\).
Oblicz długość jego najdłuższego boku.}
{\(25\, \mathrm{cm}\)}{\(12\, \mathrm{cm}\)}{\(15\, \mathrm{cm}\)}{\(20\, \mathrm{cm}\)}{\(30\, \mathrm{cm}\)}{}}
\LANGcs{
\Question{
Délky stran pravoúhlého trojúhelníka jsou tři po sobě
jdoucí členy aritmetické posloupnosti. Obvod trojúhelníka je
\(60\, \mathrm{cm}\).
Délka přepony je:}
{\(25\, \mathrm{cm}\)}{\(12\, \mathrm{cm}\)}{\(15\, \mathrm{cm}\)}{\(20\, \mathrm{cm}\)}{\(30\, \mathrm{cm}\)}{}}
\LANGsk{
\Question{
Dĺžky strán pravouhlého trojuholníka sú tri po sebe
idúce členy aritmetickej postupnosti. Obvod trojuholníka je
\(60\, \mathrm{cm}\).
Dĺžka prepony je:}
{\(25\, \mathrm{cm}\)}{\(12\, \mathrm{cm}\)}{\(15\, \mathrm{cm}\)}{\(20\, \mathrm{cm}\)}{\(30\, \mathrm{cm}\)}{}}
\LANGes{
\Question{
Los lados de un triángulo rectángulo forman tres términos consecutivos de una progresión aritmética. El perímetro
del triángulo mide \(60\, \mathrm{cm}\). ¿Cuánto mide la hipotenusa del triángulo?}
{\(25\, \mathrm{cm}\)}{\(12\, \mathrm{cm}\)}{\(15\, \mathrm{cm}\)}{\(20\, \mathrm{cm}\)}{\(30\, \mathrm{cm}\)}{}}
\End

\Data\ID{9000064802}
\Updated{2021-05-03 08:05:00}
\Level{C}
\SubArea{Arithmetic sequences}
\LANGen{
\Question{
The arithmetic sequence is given by the third term
\(a_{3} = 5\) and the common
difference \(d = 2\).
How many terms of the sequence has to be summed up to ensure that the sum is bigger
than \(300\)?}
{\(18\)}{\(10\)}{\(12\)}{\(14\)}{\(16\)}{}}
\LANGpl{
\Question{
Trzeci wyraz ciągu arytmetycznego jest równy 
\(a_{3} = 5\), a różnica
wynosi  \(d = 2\).
Ile wyrazów tego ciągu musimy do siebie dodać, aby otrzymana suma była większa 
niż \(300\)?}
{\(18\)}{\(10\)}{\(12\)}{\(14\)}{\(16\)}{}}
\LANGcs{
\Question{
V aritmetické posloupnosti je \(a_{3} = 5\),
\(d = 2\).
Kolik členů musíme sečíst, aby součet byl větší než
\(300\)?}
{\(18\)}{\(10\)}{\(12\)}{\(14\)}{\(16\)}{}}
\LANGsk{
\Question{
V aritmetickej postupnosti je \(a_{3} = 5\),
\(d = 2\).
Koľko členov musíme sčítať, aby súčet bol väčší než
\(300\)?}
{\(18\)}{\(10\)}{\(12\)}{\(14\)}{\(16\)}{}}
\LANGes{
\Question{
Una progresión aritmética viene dada por su tercer término
\(a_{3} = 5\) y la diferencia \(d = 2\).
¿Cuántos términos de la progresión tenemos que sumar para que la suma sea mayor que \(300\)?}
{\(18\)}{\(10\)}{\(12\)}{\(14\)}{\(16\)}{}}
\End

\Data\ID{9000064803}
\Updated{2020-12-28 03:12:31}
\Level{C}
\SubArea{Arithmetic sequences}
\LANGen{
\Question{
Suppose three numbers form three consecutive terms of an arithmetic sequence.  The sum of these
numbers is \(33\) and
the product is \(1\: 155\).
Find the smallest of these three numbers.}
{\(7\)}{\(9\)}{\(11\)}{\(13\)}{\(15\)}{}}
\LANGpl{
\Question{
Trzy liczby są trzema kolejnymi  składnikami ciągu arytmetycznego. Suma tych liczb jest równa \(33\),  a
iloczyn jest równy  \(1\: 155\).
Jaka jest najmniejsza z tych trzech liczb?}
{\(7\)}{\(9\)}{\(11\)}{\(13\)}{\(15\)}{}}
\LANGcs{
\Question{
Tři čísla, která tvoří aritmetickou posloupnost, mají součet
\(33\) a
součin \(1\: 155\).
Nejmenší z těchto čísel je:}
{\(7\)}{\(9\)}{\(11\)}{\(13\)}{\(15\)}{}}
\LANGsk{
\Question{
Tri čísla, ktoré tvoria aritmetickú postupnosť, majú súčet
\(33\) a
súčin \(1\: 155\).
Najmenšie z týchto troch čísel je:}
{\(7\)}{\(9\)}{\(11\)}{\(13\)}{\(15\)}{}}
\LANGes{
\Question{
Tres números forman tres términos consecutivos de una progresión aritmética. La suma de estos números es \(33\) y su producto es \(1\: 155\).
Halla el número más pequeño de estos números.}
{\(7\)}{\(9\)}{\(11\)}{\(13\)}{\(15\)}{}}
\End

\Data\ID{9000064804}
\Updated{2020-07-15 03:07:37}
\Level{C}
\SubArea{Arithmetic sequences}
\LANGen{
\Question{
An arithmetic sequence is formed by consecutive odd integers. The twelfth term is
\( 53\). Find
the sum of the first five terms.}
{\(175\)}{\(151\)}{\(163\)}{\(187\)}{\(199\)}{}}
\LANGpl{
\Question{
Ciąg arytmetyczny został utworzony z kolejnych liczb nieparzystych. Dwunasty wyraz tego ciągu to 
\(53\). Oblicz sumę pięciu początkowych wyrazów tego ciągu.}
{\(175\)}{\(151\)}{\(163\)}{\(187\)}{\(199\)}{}}
\LANGcs{
\Question{
V posloupnosti, která je tvořena po sobě jdoucími lichými čísly, platí
\(a_{12} = 53\).
Součet prvních pěti členů je:}
{\(175\)}{\(151\)}{\(163\)}{\(187\)}{\(199\)}{}}
\LANGsk{
\Question{
V postupnosti, ktorá je tvorená po sebe idúcimi nepárnymi číslami, platí
\(a_{12}=53\).
Súčet prvých piatich členov je:}
{\(175\)}{\(151\)}{\(163\)}{\(187\)}{\(199\)}{}}
\LANGes{
\Question{
Una progresión aritmética está formada por los números impares consecutivos. El término 12 es
\( 53\). Halla la suma de los cinco primeros términos.}
{\(175\)}{\(151\)}{\(163\)}{\(187\)}{\(199\)}{}}
\End

\Data\ID{9000064805}
\Updated{2020-07-15 03:07:37}
\Level{C}
\SubArea{Arithmetic sequences}
\LANGen{
\Question{
The sides of a box form three consecutive terms of an arithmetic sequence. The volume of the box is
\(665\, \mathrm{cm}^{3}\). The shortest
side is \(5\, \mathrm{cm}\).
Find the surface area of the box.}
{\(501\, \mathrm{cm}^{2}\)}{\(315\, \mathrm{cm}^{2}\)}{\(615\, \mathrm{cm}^{2}\)}{\(805\, \mathrm{cm}^{2}\)}{\(1\: 215\, \mathrm{cm}^{2}\)}{}}
\LANGpl{
\Question{
Długości boków  pudełka są trzema  wyrazami ciągu arytmetycznego. Objętość tego pudełka wynosi 
\(665\, \mathrm{cm}^{3}\), a najkrótszy bok ma długość \(5\, \mathrm{cm}\).
Oblicz pole powierzchni tego pudełka.}
{\(501\, \mathrm{cm}^{2}\)}{\(315\, \mathrm{cm}^{2}\)}{\(615\, \mathrm{cm}^{2}\)}{\(805\, \mathrm{cm}^{2}\)}{\(1\: 215\, \mathrm{cm}^{2}\)}{}}
\LANGcs{
\Question{
Délky hran kvádru tvoří tři po sobě jdoucí členy aritmetické posloupnosti. Objem kvádru je
\(665\, \mathrm{cm}^{3}\). Jeho nejkratší
hrana měří \(5\, \mathrm{cm}\).
Jeho povrch je:}
{\(501\, \mathrm{cm}^{2}\)}{\(315\, \mathrm{cm}^{2}\)}{\(615\, \mathrm{cm}^{2}\)}{\(805\, \mathrm{cm}^{2}\)}{\(1\: 215\, \mathrm{cm}^{2}\)}{}}
\LANGsk{
\Question{
Dĺžky hrán kvádra tvoria tri po sebe idúce členy aritmetickej postupnosti. Objem kvádra je
\(665\, \mathrm{cm}^{3}\). Jeho najkratšia
hrana meria \(5\, \mathrm{cm}\).
Jeho povrch je:}
{\(501\, \mathrm{cm}^{2}\)}{\(315\, \mathrm{cm}^{2}\)}{\(615\, \mathrm{cm}^{2}\)}{\(805\, \mathrm{cm}^{2}\)}{\(1\: 215\, \mathrm{cm}^{2}\)}{}}
\LANGes{
\Question{
Las aristas de un ortoedro forman tres términos consecutivos de una progresión aritmética. El volumen del ortoedro es
\(665\, \mathrm{cm}^{3}\). La arista más corta mide \(5\, \mathrm{cm}\).
Halla el área del ortoedro.}
{\(501\, \mathrm{cm}^{2}\)}{\(315\, \mathrm{cm}^{2}\)}{\(615\, \mathrm{cm}^{2}\)}{\(805\, \mathrm{cm}^{2}\)}{\(1\: 215\, \mathrm{cm}^{2}\)}{}}
\End

\Data\ID{9000064807}
\Updated{2020-07-15 03:07:37}
\Level{C}
\SubArea{Arithmetic sequences}
\LANGen{
\Question{
Find the sum of the integers which satisfy the following inequality. 
\[ 
 x^{2} - 8x - 153\leq 0
\]}
{\(108\)}{\(162\)}{\(91\)}{\(78\)}{\(56\)}{}}
\LANGpl{
\Question{
Oblicz sumę wszystkich liczb całkowitych, które spełniają następującą nierówność. 
\[ 
 x^{2} - 8x - 153\leq 0
\]}
{\(108\)}{\(162\)}{\(91\)}{\(78\)}{\(56\)}{}}
\LANGcs{
\Question{
Určete součet všech celých čísel, které vyhovují nerovnici
\(x^{2} - 8x - 153\leq 0\).}
{\(108\)}{\(162\)}{\(91\)}{\(78\)}{\(56\)}{}}
\LANGsk{
\Question{
Určte súčet všetkých celých čísel, ktoré vyhovujú nerovnici
\(x^{2} - 8x - 153\leq 0\).}
{\(108\)}{\(162\)}{\(91\)}{\(78\)}{\(56\)}{}}
\LANGes{
\Question{
Halla la suma de todos los números enteros que satisfacen la siguiente inecuación: 
\[ 
 x^{2} - 8x - 153\leq 0
\]}
{\(108\)}{\(162\)}{\(91\)}{\(78\)}{\(56\)}{}}
\End

\Data\ID{9000064808}
\Updated{2021-05-03 08:05:13}
\Level{C}
\SubArea{Arithmetic sequences}
\LANGen{
\Question{
The sum of the first eight terms of an arithmetic sequence is
\(44\).
The sum of the next four terms is bigger than this value by
\(50\). Find the
thirteenth term \(a_{13}\).}
{\(31\)}{\(22\)}{\(25\)}{\(28\)}{\(34\)}{}}
\LANGpl{
\Question{
Suma ośmiu początkowych wyrazów ciągu arytmetycznego wynosi  \(44\).
Suma czterech kolejnych wyrazów tego ciągu jest większa niż powyższa wartość o
\(50\). Jaki jest trzynasty wyraz tego ciągu \(a_{13}\)?}
{\(31\)}{\(22\)}{\(25\)}{\(28\)}{\(34\)}{}}
\LANGcs{
\Question{
Součet prvních osmi členů aritmetické posloupnosti je
\(44\). Součet následujících
čtyř členů je o \(50\) 
větší. Třináctý člen posloupnosti je:}
{\(31\)}{\(22\)}{\(25\)}{\(28\)}{\(34\)}{}}
\LANGsk{
\Question{
Súčet prvých ôsmich členov aritmetickej postupnosti je
\(44\). Súčet nasledujúcich
štyroch členov je o \(50\) 
väčší. Trinásty člen postupnosti je:}
{\(31\)}{\(22\)}{\(25\)}{\(28\)}{\(34\)}{}}
\LANGes{
\Question{
La suma de los ocho primeros términos de una progresión aritmética es
\(44\).
La suma de los siguiente cuatro términos es mayor que dicho valor en
\(50\) unidades. Halla el término trece \(a_{13}\).}
{\(31\)}{\(22\)}{\(25\)}{\(28\)}{\(34\)}{}}
\End
